
\documentclass{article}
\usepackage{colortbl}
\usepackage{makecell}
\usepackage{multirow}
\usepackage{supertabular}

\begin{document}

\newcounter{utterance}

\centering \large Interaction Transcript for game `clean\_up', experiment `0\_easy\_5obj\_mt', episode 0 with nemotron{-}3{-}ultra.
\vspace{24pt}

{ \footnotesize  \setcounter{utterance}{1}
\setlength{\tabcolsep}{0pt}
\begin{supertabular}{c@{$\;$}|p{.15\linewidth}@{}p{.15\linewidth}p{.15\linewidth}p{.15\linewidth}p{.15\linewidth}p{.15\linewidth}}
    \# & $\;$A & \multicolumn{4}{c}{Game Master} & $\;\:$B\\
    \hline

    \theutterance \stepcounter{utterance}  
    & & \multicolumn{4}{p{0.6\linewidth}}{
        \cellcolor[rgb]{0.9,0.9,0.9}{
            \makecell[{{p{\linewidth}}}]{
                \texttt{\tiny{[P1$\langle$GM]}}
                \texttt{Jiena l{-}kap tal{-}logħba tiegħek, u inti qed tilgħab logħba kollaborattiva bil{-}grilja li ġejja bħala l{-}bord tal{-}logħba:} \\
\\ 
\texttt{```} \\
\texttt{ 1234567} \\
\texttt{╔═╤═══╤═╗} \\
\texttt{║◌│◌◌◌│◌║ 1} \\
\texttt{╟─┼──┬┘◌║ 2} \\
\texttt{║◌│◌◌│◌◌║ 3} \\
\texttt{╟─┴──┘H◌║ 4} \\
\texttt{║◌◌◌◌◌I◌║ 5} \\
\texttt{║C◌◌◌◌◌◌║ 6} \\
\texttt{║◌W◌◌T◌◌║ 7} \\
\texttt{╚═══════╝} \\
\\ 
\texttt{```} \\
\\ 
\texttt{* It{-}tarf ta' fuq juri l{-}koordinati x li jiżdiedu lejn il{-}lemin, u t{-}tarf tal{-}lemin il{-}koordinati y li jiżdiedu 'l isfel.} \\
\texttt{* L{-}oġġetti li ġejjin huma mqiegħda b'mod każwali fuq il{-}grilja tiegħek: 'W', 'I', 'T', 'C', 'H'.} \\
\\ 
\\ 
\texttt{Il{-}plejer l{-}ieħor jara varjazzjoni tal{-}bord tal{-}logħba, fejn l{-}oġġetti huma mqiegħda f'postijiet każwali differenti. Int ma tistax tara l{-}bord tal{-}plejer l{-}ieħor, u lanqas hu/hi ma jista'/tista' jara tiegħek.} \\
\\ 
\\ 
\texttt{**Għan:**} \\
\\ 
\texttt{Iż{-}żewġ plejers iridu jċaqalqu l{-}oġġetti fuq l{-}isfond rispettiv tagħhom sabiex oġġetti identiċi jispiċċaw fl{-}istess koordinati. Trid tikkomunika mal{-}plejer l{-}ieħor biex taqblu fuq konfigurazzjoni komuni tal{-}għan.} \\
\\ 
\\ 
\texttt{**Regoli:**} \\
\\ 
\texttt{* F'kull dawra, tista' tibgħat eżattament wieħed miż{-}żewġ kmandi li ġejjin:} \\
\texttt{ 1. `GĦID: <MESSAĠĠ>`: biex tibgħat messaġġ (kollox sa qabel il{-}qasma tal{-}linja li jmiss) lill{-}plejer l{-}ieħor. Jiena ngħaddih lis{-}sieħeb/sieħba tiegħek.} \\
\texttt{2. `ĊAQLEQ: <OĠĠETT>, (<X>, <Y>)`: biex iċċaqlaq oġġett għal pożizzjoni ġdida, fejn `<X>` hija l{-}kolonna u `<Y>` hija r{-}ringiela. Ngħidlek jekk iċ{-}ċaqliq tiegħek kienx ta' suċċess jew le.} \\
\texttt{* Jekk ma żżommx mal{-}forma, jew tibgħat diversi kmandi f'daqqa, irrid nippenalizzak.} \\
\texttt{* Jekk iż{-}żewġ plejers jakkumulaw aktar minn 12 penalitajiet, it{-}tnejn titlef il{-}logħba.} \\
\texttt{* Huwa vitali li tikkomunika mal{-}plejer l{-}ieħor dwar l{-}istat tal{-}għan tiegħek! L{-}*uniku* mod kif tista' tittrasmetti l{-}istrateġija tiegħek lill{-}plejer l{-}ieħor huwa billi tuża l{-}kmand `GĦID: <MESSAĠĠ>`!} \\
\\ 
\\ 
\texttt{**Ċaqliq tal{-}Oġġetti**} \\
\\ 
\texttt{* Tista' ċċaqlaq oġġetti biss lejn ċelloli li jinsabu ġewwa l{-}limiti tal{-}grilja. Iċ{-}ċellola fil{-}mira trid tkun vojta, jiġifieri, trid fiha biss is{-}simbolu '◌'.} \\
\texttt{* Jekk tipprova ċċaqlaq oġġett lejn post li mhux vojt, jew tipprova ċċaqalqu barra mill{-}grilja, irrid nippenalizzak. Ikollok prova oħra.} \\
\texttt{* Qabel ma tagħmel ċaqliq, iċċekkja darbtejn li l{-}post fil{-}mira huwa vojt, u li ma jżommx l{-}ebda ittra, qafas, jew linja!} \\
\\ 
\\ 
\texttt{**Tmiem il{-}Logħba**} \\
\\ 
\texttt{Jekk taħseb li wasalt għall{-}għan li tallinjaw l{-}oġġetti kollha, tista' titlob lill{-}plejer l{-}ieħor biex itemm il{-}logħba billi tibgħat `GĦID: lestejt?`. Jekk il{-}plejer l{-}ieħor jitlobek biex itemm il{-}logħba, u inti twieġeb b'`GĦID: lestejt!`, il{-}logħba tintemm.} \\
\\ 
\texttt{Iż{-}żewġ plejers jirbħu jekk il{-}logħba tispiċċa fi żmien 20 rawnd, fejn rawnd wieħed huwa definit bħala żewġ plejers li kull wieħed jibgħat kmand validu.} \\
\\ 
\texttt{**Punteġġ:**} \\
\\ 
\texttt{Iktar ma l{-}oġġetti identiċi jkunu qrib xulxin fiż{-}żewġ bordijiet tal{-}logħba, iktar taqla' punti. Il{-}penalitajiet inaqqsu l{-}punti tiegħek. Tista' tgħaddi r{-}rekord?} \\
\\ 
\\ 
\texttt{Jekk jogħġbok ibgħat messaġġ lill{-}plejer l{-}ieħor biex tibda l{-}logħba!} \\
            }
        }
    }
    & \\ \\

    \theutterance \stepcounter{utterance}  
    & \multicolumn{4}{p{0.6\linewidth}}{
        \cellcolor[rgb]{0.9,0.9,0.9}{
            \makecell[{{p{\linewidth}}}]{
                \texttt{\tiny{[P1$\rangle$GM]}}
                \texttt{GĦID: Merħba! Nibda nara l{-}oġġetti tiegħi: W f'(2,7), I f'(5,5), T f'(4,7), C f'(1,6), H f'(7,4). Fejn qegħdin l{-}oġġetti tiegħek?} \\
            }
        }
    }
    & & \\ \\

    \theutterance \stepcounter{utterance}  
    & & & \multicolumn{2}{p{0.3\linewidth}}{
        \cellcolor[rgb]{0.9,0.9,0.9}{
            \makecell[{{p{\linewidth}}}]{
                \texttt{\tiny{[GM$|$GM]}}
                \texttt{GĦID: Merħba! Nibda nara l{-}oġġetti tiegħi: W f'(2,7), I f'(5,5), T f'(4,7), C f'(1,6), H f'(7,4). Fejn qegħdin l{-}oġġetti tiegħek?} \\
            }
        }
    }
    & & \\ \\

    \theutterance \stepcounter{utterance}  
    & & \multicolumn{4}{p{0.6\linewidth}}{
        \cellcolor[rgb]{0.9,0.9,0.9}{
            \makecell[{{p{\linewidth}}}]{
                \texttt{\tiny{[GM$\rangle$P2]}}
                \texttt{Jiena l{-}kap tal{-}logħba tiegħek, u inti qed tilgħab logħba kollaborattiva bil{-}grilja li ġejja bħala l{-}bord tal{-}logħba:} \\
\\ 
\texttt{```} \\
\texttt{ 1234567} \\
\texttt{╔═╤═══╤═╗} \\
\texttt{║◌│◌◌◌│H║ 1} \\
\texttt{╟─┼──┬┘◌║ 2} \\
\texttt{║◌│◌◌│◌C║ 3} \\
\texttt{╟─┴──┘◌T║ 4} \\
\texttt{║◌◌◌◌◌◌◌║ 5} \\
\texttt{║◌◌◌I◌◌◌║ 6} \\
\texttt{║◌◌◌◌W◌◌║ 7} \\
\texttt{╚═══════╝} \\
\\ 
\texttt{```} \\
\\ 
\texttt{* It{-}tarf ta' fuq juri l{-}koordinati x li jiżdiedu lejn il{-}lemin, u t{-}tarf tal{-}lemin il{-}koordinati y li jiżdiedu 'l isfel.} \\
\texttt{* L{-}oġġetti li ġejjin huma mqiegħda b'mod każwali fuq il{-}grilja tiegħek: 'W', 'I', 'T', 'C', 'H'.} \\
\\ 
\\ 
\texttt{Il{-}plejer l{-}ieħor jara varjazzjoni tal{-}bord tal{-}logħba, fejn l{-}oġġetti huma mqiegħda f'postijiet każwali differenti. Int ma tistax tara l{-}bord tal{-}plejer l{-}ieħor, u lanqas hu/hi ma jista'/tista' jara tiegħek.} \\
\\ 
\\ 
\texttt{**Għan:**} \\
\\ 
\texttt{Iż{-}żewġ plejers iridu jċaqalqu l{-}oġġetti fuq l{-}isfond rispettiv tagħhom sabiex oġġetti identiċi jispiċċaw fl{-}istess koordinati. Trid tikkomunika mal{-}plejer l{-}ieħor biex taqblu fuq konfigurazzjoni komuni tal{-}għan.} \\
\\ 
\\ 
\texttt{**Regoli:**} \\
\\ 
\texttt{* F'kull dawra, tista' tibgħat eżattament wieħed miż{-}żewġ kmandi li ġejjin:} \\
\texttt{ 1. `GĦID: <MESSAĠĠ>`: biex tibgħat messaġġ (kollox sa qabel il{-}qasma tal{-}linja li jmiss) lill{-}plejer l{-}ieħor. Jiena ngħaddih lis{-}sieħeb/sieħba tiegħek.} \\
\texttt{2. `ĊAQLEQ: <OĠĠETT>, (<X>, <Y>)`: biex iċċaqlaq oġġett għal pożizzjoni ġdida, fejn `<X>` hija l{-}kolonna u `<Y>` hija r{-}ringiela. Ngħidlek jekk iċ{-}ċaqliq tiegħek kienx ta' suċċess jew le.} \\
\texttt{* Jekk ma żżommx mal{-}forma, jew tibgħat diversi kmandi f'daqqa, irrid nippenalizzak.} \\
\texttt{* Jekk iż{-}żewġ plejers jakkumulaw aktar minn 12 penalitajiet, it{-}tnejn titlef il{-}logħba.} \\
\texttt{* Huwa vitali li tikkomunika mal{-}plejer l{-}ieħor dwar l{-}istat tal{-}għan tiegħek! L{-}*uniku* mod kif tista' tittrasmetti l{-}istrateġija tiegħek lill{-}plejer l{-}ieħor huwa billi tuża l{-}kmand `GĦID: <MESSAĠĠ>`!} \\
\\ 
\\ 
\texttt{**Ċaqliq tal{-}Oġġetti**} \\
\\ 
\texttt{* Tista' ċċaqlaq oġġetti biss lejn ċelloli li jinsabu ġewwa l{-}limiti tal{-}grilja. Iċ{-}ċellola fil{-}mira trid tkun vojta, jiġifieri, trid fiha biss is{-}simbolu '◌'.} \\
\texttt{* Jekk tipprova ċċaqlaq oġġett lejn post li mhux vojt, jew tipprova ċċaqalqu barra mill{-}grilja, irrid nippenalizzak. Ikollok prova oħra.} \\
\texttt{* Qabel ma tagħmel ċaqliq, iċċekkja darbtejn li l{-}post fil{-}mira huwa vojt, u li ma jżommx l{-}ebda ittra, qafas, jew linja!} \\
\\ 
\\ 
\texttt{**Tmiem il{-}Logħba**} \\
\\ 
\texttt{Jekk taħseb li wasalt għall{-}għan li tallinjaw l{-}oġġetti kollha, tista' titlob lill{-}plejer l{-}ieħor biex itemm il{-}logħba billi tibgħat `GĦID: lestejt?`. Jekk il{-}plejer l{-}ieħor jitlobek biex itemm il{-}logħba, u inti twieġeb b'`GĦID: lestejt!`, il{-}logħba tintemm.} \\
\\ 
\texttt{Iż{-}żewġ plejers jirbħu jekk il{-}logħba tispiċċa fi żmien 20 rawnd, fejn rawnd wieħed huwa definit bħala żewġ plejers li kull wieħed jibgħat kmand validu.} \\
\\ 
\texttt{**Punteġġ:**} \\
\\ 
\texttt{Iktar ma l{-}oġġetti identiċi jkunu qrib xulxin fiż{-}żewġ bordijiet tal{-}logħba, iktar taqla' punti. Il{-}penalitajiet inaqqsu l{-}punti tiegħek. Tista' tgħaddi r{-}rekord?} \\
\\ 
\texttt{Il{-}plejer l{-}ieħor beda l{-}logħba billi bagħat dan il{-}messaġġ:} \\
\\ 
\texttt{"Merħba! Nibda nara l{-}oġġetti tiegħi: W f'(2,7), I f'(5,5), T f'(4,7), C f'(1,6), H f'(7,4). Fejn qegħdin l{-}oġġetti tiegħek?"} \\
\\ 
\texttt{X'inhu l{-}ewwel kmand tiegħek?} \\
            }
        }
    }
    & \\ \\

    \theutterance \stepcounter{utterance}  
    & & & \multicolumn{4}{p{0.6\linewidth}}{
        \cellcolor[rgb]{0.9,0.9,0.9}{
            \makecell[{{p{\linewidth}}}]{
                \texttt{\tiny{[GM$\langle$P2]}}
                \texttt{GĦID: L{-}oġġetti tiegħi huma: W f'(7,7), I f'(4,6), T f'(4,4), C f'(7,3), H f'(7,1). Nitgħallem li bordijiet differenti għandhom pożizzjonijiet differenti għall{-}oġġetti. Nitgħallmu qiegħ l{-}għan komuni: nistgħu niżommu l{-}oġġetti fil{-}bord: (1,1)=W, (2,1)=I, (3,1)=T, (4,1)=C, (5,1)=H? Jekk inti tagħmel dan, nistgħu nispiċċaw l{-}oġġetti. Jekk inti tagħmel li huwa possibbli, għidli "tajjeb".} \\
            }
        }
    }
     \\ \\

    \theutterance \stepcounter{utterance}  
    & & & \multicolumn{2}{p{0.3\linewidth}}{
        \cellcolor[rgb]{0.9,0.9,0.9}{
            \makecell[{{p{\linewidth}}}]{
                \texttt{\tiny{[GM$|$GM]}}
                \texttt{GĦID: L{-}oġġetti tiegħi huma: W f'(7,7), I f'(4,6), T f'(4,4), C f'(7,3), H f'(7,1). Nitgħallem li bordijiet differenti għandhom pożizzjonijiet differenti għall{-}oġġetti. Nitgħallmu qiegħ l{-}għan komuni: nistgħu niżommu l{-}oġġetti fil{-}bord: (1,1)=W, (2,1)=I, (3,1)=T, (4,1)=C, (5,1)=H? Jekk inti tagħmel dan, nistgħu nispiċċaw l{-}oġġetti. Jekk inti tagħmel li huwa possibbli, għidli "tajjeb".} \\
            }
        }
    }
    & & \\ \\

    \theutterance \stepcounter{utterance}  
    & & \multicolumn{4}{p{0.6\linewidth}}{
        \cellcolor[rgb]{0.9,0.9,0.9}{
            \makecell[{{p{\linewidth}}}]{
                \texttt{\tiny{[P1$\langle$GM]}}
                \texttt{Il{-}messaġġ tiegħek ġie mgħoddi lill{-}plejer l{-}ieħor.} \\
\\ 
\\ 
\texttt{Bħalissa qed tilagħbu r{-}round 1 minn massimu ta' 20.} \\
\texttt{Intom flimkien akkumulajtu 0 minn 12 penali.} \\
\texttt{Il{-}plejer l{-}ieħor bagħat dan il{-}messaġġ:} \\
\\ 
\\ 
\texttt{"L{-}oġġetti tiegħi huma: W f'(7,7), I f'(4,6), T f'(4,4), C f'(7,3), H f'(7,1). Nitgħallem li bordijiet differenti għandhom pożizzjonijiet differenti għall{-}oġġetti. Nitgħallmu qiegħ l{-}għan komuni: nistgħu niżommu l{-}oġġetti fil{-}bord: (1,1)=W, (2,1)=I, (3,1)=T, (4,1)=C, (5,1)=H? Jekk inti tagħmel dan, nistgħu nispiċċaw l{-}oġġetti. Jekk inti tagħmel li huwa possibbli, għidli "tajjeb"."} \\
\\ 
\\ 
\texttt{X'inhu l{-}kmand li jmiss tiegħek?} \\
            }
        }
    }
    & \\ \\

    \theutterance \stepcounter{utterance}  
    & \multicolumn{4}{p{0.6\linewidth}}{
        \cellcolor[rgb]{0.9,0.9,0.9}{
            \makecell[{{p{\linewidth}}}]{
                \texttt{\tiny{[P1$\rangle$GM]}}
                \texttt{GĦID: Tajjeb, niqqasem. Sef inħareg l{-}oġġetti f'riga waħda fuq ir{-}ringiela 1. Inbidu m'W għal (1,1).} \\
            }
        }
    }
    & & \\ \\

    \theutterance \stepcounter{utterance}  
    & & & \multicolumn{2}{p{0.3\linewidth}}{
        \cellcolor[rgb]{0.9,0.9,0.9}{
            \makecell[{{p{\linewidth}}}]{
                \texttt{\tiny{[GM$|$GM]}}
                \texttt{GĦID: Tajjeb, niqqasem. Sef inħareg l{-}oġġetti f'riga waħda fuq ir{-}ringiela 1. Inbidu m'W għal (1,1).} \\
            }
        }
    }
    & & \\ \\

    \theutterance \stepcounter{utterance}  
    & & \multicolumn{4}{p{0.6\linewidth}}{
        \cellcolor[rgb]{0.9,0.9,0.9}{
            \makecell[{{p{\linewidth}}}]{
                \texttt{\tiny{[GM$\rangle$P2]}}
                \texttt{Il{-}messaġġ tiegħek ġie mgħoddi lill{-}plejer l{-}ieħor.} \\
\\ 
\\ 
\texttt{Bħalissa qed tilagħbu r{-}round 2 minn massimu ta' 20.} \\
\texttt{Intom flimkien akkumulajtu 0 minn 12 penali.} \\
\texttt{Il{-}plejer l{-}ieħor bagħat dan il{-}messaġġ:} \\
\\ 
\\ 
\texttt{"Tajjeb, niqqasem. Sef inħareg l{-}oġġetti f'riga waħda fuq ir{-}ringiela 1. Inbidu m'W għal (1,1)."} \\
\\ 
\\ 
\texttt{X'inhu l{-}kmand li jmiss tiegħek?} \\
            }
        }
    }
    & \\ \\

    \theutterance \stepcounter{utterance}  
    & & & \multicolumn{4}{p{0.6\linewidth}}{
        \cellcolor[rgb]{0.9,0.9,0.9}{
            \makecell[{{p{\linewidth}}}]{
                \texttt{\tiny{[GM$\langle$P2]}}
                \texttt{ĊAQLEQ: W, (1, 1)} \\
            }
        }
    }
     \\ \\

    \theutterance \stepcounter{utterance}  
    & & & \multicolumn{2}{p{0.3\linewidth}}{
        \cellcolor[rgb]{0.9,0.9,0.9}{
            \makecell[{{p{\linewidth}}}]{
                \texttt{\tiny{[GM$|$GM]}}
                \texttt{ĊAQLEQ: W, (1, 1)} \\
            }
        }
    }
    & & \\ \\

    \theutterance \stepcounter{utterance}  
    & & & \multicolumn{2}{p{0.3\linewidth}}{
        \cellcolor[rgb]{0.9,0.9,0.9}{
            \makecell[{{p{\linewidth}}}]{
                \texttt{\tiny{[GM$|$GM]}}
                \texttt{W moved to (1, 1)} \\
            }
        }
    }
    & & \\ \\

    \theutterance \stepcounter{utterance}  
    & & \multicolumn{4}{p{0.6\linewidth}}{
        \cellcolor[rgb]{0.9,0.9,0.9}{
            \makecell[{{p{\linewidth}}}]{
                \texttt{\tiny{[P1$\langle$GM]}}
                \texttt{Il{-}messaġġ tiegħek ġie mgħoddi lill{-}plejer l{-}ieħor.} \\
\\ 
\\ 
\texttt{Bħalissa qed tilagħbu r{-}round 2 minn massimu ta' 20.} \\
\texttt{Intom flimkien akkumulajtu 0 minn 12 penali.} \\
\texttt{Il{-}plejer l{-}ieħor mexxa oġġett fuq il{-}grilja tiegħu.} \\
\\ 
\\ 
\texttt{X'inhu l{-}kmand li jmiss tiegħek?} \\
            }
        }
    }
    & \\ \\

    \theutterance \stepcounter{utterance}  
    & \multicolumn{4}{p{0.6\linewidth}}{
        \cellcolor[rgb]{0.9,0.9,0.9}{
            \makecell[{{p{\linewidth}}}]{
                \texttt{\tiny{[P1$\rangle$GM]}}
                \texttt{ĊAQLEQ: W, (1, 1)} \\
            }
        }
    }
    & & \\ \\

    \theutterance \stepcounter{utterance}  
    & & & \multicolumn{2}{p{0.3\linewidth}}{
        \cellcolor[rgb]{0.9,0.9,0.9}{
            \makecell[{{p{\linewidth}}}]{
                \texttt{\tiny{[GM$|$GM]}}
                \texttt{ĊAQLEQ: W, (1, 1)} \\
            }
        }
    }
    & & \\ \\

    \theutterance \stepcounter{utterance}  
    & & & \multicolumn{2}{p{0.3\linewidth}}{
        \cellcolor[rgb]{0.9,0.9,0.9}{
            \makecell[{{p{\linewidth}}}]{
                \texttt{\tiny{[GM$|$GM]}}
                \texttt{W moved to (1, 1)} \\
            }
        }
    }
    & & \\ \\

    \theutterance \stepcounter{utterance}  
    & & \multicolumn{4}{p{0.6\linewidth}}{
        \cellcolor[rgb]{0.9,0.9,0.9}{
            \makecell[{{p{\linewidth}}}]{
                \texttt{\tiny{[GM$\rangle$P2]}}
                \texttt{Mexxejt 'W' għal (1,1) b'suċċess. Il{-}grilja aġġornata tiegħek tidher hekk:} \\
\texttt{```} \\
\texttt{ 1234567} \\
\texttt{╔═╤═══╤═╗} \\
\texttt{║W│◌◌◌│H║ 1} \\
\texttt{╟─┼──┬┘◌║ 2} \\
\texttt{║◌│◌◌│◌C║ 3} \\
\texttt{╟─┴──┘◌T║ 4} \\
\texttt{║◌◌◌◌◌◌◌║ 5} \\
\texttt{║◌◌◌I◌◌◌║ 6} \\
\texttt{║◌◌◌◌◌◌◌║ 7} \\
\texttt{╚═══════╝} \\
\texttt{```} \\
\texttt{Bħalissa qed tilagħbu r{-}round 3 minn massimu ta' 20.} \\
\texttt{Intom flimkien akkumulajtu 0 minn 12 penali.} \\
\texttt{Il{-}plejer l{-}ieħor mexxa oġġett fuq il{-}grilja tiegħu.} \\
\\ 
\\ 
\texttt{X'inhu l{-}kmand li jmiss tiegħek?} \\
            }
        }
    }
    & \\ \\

    \theutterance \stepcounter{utterance}  
    & & & \multicolumn{4}{p{0.6\linewidth}}{
        \cellcolor[rgb]{0.9,0.9,0.9}{
            \makecell[{{p{\linewidth}}}]{
                \texttt{\tiny{[GM$\langle$P2]}}
                \texttt{ĊAQLEQ: I, (2, 1)} \\
            }
        }
    }
     \\ \\

    \theutterance \stepcounter{utterance}  
    & & & \multicolumn{2}{p{0.3\linewidth}}{
        \cellcolor[rgb]{0.9,0.9,0.9}{
            \makecell[{{p{\linewidth}}}]{
                \texttt{\tiny{[GM$|$GM]}}
                \texttt{ĊAQLEQ: I, (2, 1)} \\
            }
        }
    }
    & & \\ \\

    \theutterance \stepcounter{utterance}  
    & & & \multicolumn{2}{p{0.3\linewidth}}{
        \cellcolor[rgb]{0.9,0.9,0.9}{
            \makecell[{{p{\linewidth}}}]{
                \texttt{\tiny{[GM$|$GM]}}
                \texttt{Penali: (2,1) mhijiex vojta, iżda fiha '│'.} \\
\\ 
\texttt{Intom flimkien akkumulajtu 1 minn 12 penali.} \\
\texttt{Jekk jogħġbok erġa' pprova!} \\
            }
        }
    }
    & & \\ \\

    \theutterance \stepcounter{utterance}  
    & & \multicolumn{4}{p{0.6\linewidth}}{
        \cellcolor[rgb]{0.9,0.9,0.9}{
            \makecell[{{p{\linewidth}}}]{
                \texttt{\tiny{[GM$\rangle$P2]}}
                \texttt{Penali: (2,1) mhijiex vojta, iżda fiha '│'.} \\
\\ 
\texttt{Intom flimkien akkumulajtu 1 minn 12 penali.} \\
\texttt{Jekk jogħġbok erġa' pprova!} \\
            }
        }
    }
    & \\ \\

    \theutterance \stepcounter{utterance}  
    & & & \multicolumn{4}{p{0.6\linewidth}}{
        \cellcolor[rgb]{0.9,0.9,0.9}{
            \makecell[{{p{\linewidth}}}]{
                \texttt{\tiny{[GM$\langle$P2]}}
                \texttt{ĊAQLEQ: I, (3, 1)} \\
            }
        }
    }
     \\ \\

    \theutterance \stepcounter{utterance}  
    & & & \multicolumn{2}{p{0.3\linewidth}}{
        \cellcolor[rgb]{0.9,0.9,0.9}{
            \makecell[{{p{\linewidth}}}]{
                \texttt{\tiny{[GM$|$GM]}}
                \texttt{ĊAQLEQ: I, (3, 1)} \\
            }
        }
    }
    & & \\ \\

    \theutterance \stepcounter{utterance}  
    & & & \multicolumn{2}{p{0.3\linewidth}}{
        \cellcolor[rgb]{0.9,0.9,0.9}{
            \makecell[{{p{\linewidth}}}]{
                \texttt{\tiny{[GM$|$GM]}}
                \texttt{I moved to (3, 1)} \\
            }
        }
    }
    & & \\ \\

    \theutterance \stepcounter{utterance}  
    & & \multicolumn{4}{p{0.6\linewidth}}{
        \cellcolor[rgb]{0.9,0.9,0.9}{
            \makecell[{{p{\linewidth}}}]{
                \texttt{\tiny{[P1$\langle$GM]}}
                \texttt{Mexxejt 'W' għal (1,1) b'suċċess. Il{-}grilja aġġornata tiegħek tidher hekk:} \\
\texttt{```} \\
\texttt{ 1234567} \\
\texttt{╔═╤═══╤═╗} \\
\texttt{║W│◌◌◌│◌║ 1} \\
\texttt{╟─┼──┬┘◌║ 2} \\
\texttt{║◌│◌◌│◌◌║ 3} \\
\texttt{╟─┴──┘H◌║ 4} \\
\texttt{║◌◌◌◌◌I◌║ 5} \\
\texttt{║C◌◌◌◌◌◌║ 6} \\
\texttt{║◌◌◌◌T◌◌║ 7} \\
\texttt{╚═══════╝} \\
\texttt{```} \\
\texttt{Bħalissa qed tilagħbu r{-}round 3 minn massimu ta' 20.} \\
\texttt{Intom flimkien akkumulajtu 1 minn 12 penali.} \\
\texttt{Il{-}plejer l{-}ieħor mexxa oġġett fuq il{-}grilja tiegħu.} \\
\\ 
\\ 
\texttt{X'inhu l{-}kmand li jmiss tiegħek?} \\
            }
        }
    }
    & \\ \\

    \theutterance \stepcounter{utterance}  
    & \multicolumn{4}{p{0.6\linewidth}}{
        \cellcolor[rgb]{0.9,0.9,0.9}{
            \makecell[{{p{\linewidth}}}]{
                \texttt{\tiny{[P1$\rangle$GM]}}
                \texttt{ĊAQLEQ: I, (2, 1)} \\
            }
        }
    }
    & & \\ \\

    \theutterance \stepcounter{utterance}  
    & & & \multicolumn{2}{p{0.3\linewidth}}{
        \cellcolor[rgb]{0.9,0.9,0.9}{
            \makecell[{{p{\linewidth}}}]{
                \texttt{\tiny{[GM$|$GM]}}
                \texttt{ĊAQLEQ: I, (2, 1)} \\
            }
        }
    }
    & & \\ \\

    \theutterance \stepcounter{utterance}  
    & & & \multicolumn{2}{p{0.3\linewidth}}{
        \cellcolor[rgb]{0.9,0.9,0.9}{
            \makecell[{{p{\linewidth}}}]{
                \texttt{\tiny{[GM$|$GM]}}
                \texttt{Penali: (2,1) mhijiex vojta, iżda fiha '│'.} \\
\\ 
\texttt{Intom flimkien akkumulajtu 2 minn 12 penali.} \\
\texttt{Jekk jogħġbok erġa' pprova!} \\
            }
        }
    }
    & & \\ \\

    \theutterance \stepcounter{utterance}  
    & & \multicolumn{4}{p{0.6\linewidth}}{
        \cellcolor[rgb]{0.9,0.9,0.9}{
            \makecell[{{p{\linewidth}}}]{
                \texttt{\tiny{[P1$\langle$GM]}}
                \texttt{Penali: (2,1) mhijiex vojta, iżda fiha '│'.} \\
\\ 
\texttt{Intom flimkien akkumulajtu 2 minn 12 penali.} \\
\texttt{Jekk jogħġbok erġa' pprova!} \\
            }
        }
    }
    & \\ \\

    \theutterance \stepcounter{utterance}  
    & \multicolumn{4}{p{0.6\linewidth}}{
        \cellcolor[rgb]{0.9,0.9,0.9}{
            \makecell[{{p{\linewidth}}}]{
                \texttt{\tiny{[P1$\rangle$GM]}}
                \texttt{ĊAQLEQ: I, (7, 1)} \\
            }
        }
    }
    & & \\ \\

    \theutterance \stepcounter{utterance}  
    & & & \multicolumn{2}{p{0.3\linewidth}}{
        \cellcolor[rgb]{0.9,0.9,0.9}{
            \makecell[{{p{\linewidth}}}]{
                \texttt{\tiny{[GM$|$GM]}}
                \texttt{ĊAQLEQ: I, (7, 1)} \\
            }
        }
    }
    & & \\ \\

    \theutterance \stepcounter{utterance}  
    & & & \multicolumn{2}{p{0.3\linewidth}}{
        \cellcolor[rgb]{0.9,0.9,0.9}{
            \makecell[{{p{\linewidth}}}]{
                \texttt{\tiny{[GM$|$GM]}}
                \texttt{I moved to (7, 1)} \\
            }
        }
    }
    & & \\ \\

    \theutterance \stepcounter{utterance}  
    & & \multicolumn{4}{p{0.6\linewidth}}{
        \cellcolor[rgb]{0.9,0.9,0.9}{
            \makecell[{{p{\linewidth}}}]{
                \texttt{\tiny{[GM$\rangle$P2]}}
                \texttt{Mexxejt 'I' għal (3,1) b'suċċess. Il{-}grilja aġġornata tiegħek tidher hekk:} \\
\texttt{```} \\
\texttt{ 1234567} \\
\texttt{╔═╤═══╤═╗} \\
\texttt{║W│I◌◌│H║ 1} \\
\texttt{╟─┼──┬┘◌║ 2} \\
\texttt{║◌│◌◌│◌C║ 3} \\
\texttt{╟─┴──┘◌T║ 4} \\
\texttt{║◌◌◌◌◌◌◌║ 5} \\
\texttt{║◌◌◌◌◌◌◌║ 6} \\
\texttt{║◌◌◌◌◌◌◌║ 7} \\
\texttt{╚═══════╝} \\
\texttt{```} \\
\texttt{Bħalissa qed tilagħbu r{-}round 4 minn massimu ta' 20.} \\
\texttt{Intom flimkien akkumulajtu 2 minn 12 penali.} \\
\texttt{Il{-}plejer l{-}ieħor mexxa oġġett fuq il{-}grilja tiegħu.} \\
\\ 
\\ 
\texttt{X'inhu l{-}kmand li jmiss tiegħek?} \\
            }
        }
    }
    & \\ \\

    \theutterance \stepcounter{utterance}  
    & & & \multicolumn{4}{p{0.6\linewidth}}{
        \cellcolor[rgb]{0.9,0.9,0.9}{
            \makecell[{{p{\linewidth}}}]{
                \texttt{\tiny{[GM$\langle$P2]}}
                \texttt{ĊAQLEQ: T, (4, 1)} \\
            }
        }
    }
     \\ \\

    \theutterance \stepcounter{utterance}  
    & & & \multicolumn{2}{p{0.3\linewidth}}{
        \cellcolor[rgb]{0.9,0.9,0.9}{
            \makecell[{{p{\linewidth}}}]{
                \texttt{\tiny{[GM$|$GM]}}
                \texttt{ĊAQLEQ: T, (4, 1)} \\
            }
        }
    }
    & & \\ \\

    \theutterance \stepcounter{utterance}  
    & & & \multicolumn{2}{p{0.3\linewidth}}{
        \cellcolor[rgb]{0.9,0.9,0.9}{
            \makecell[{{p{\linewidth}}}]{
                \texttt{\tiny{[GM$|$GM]}}
                \texttt{T moved to (4, 1)} \\
            }
        }
    }
    & & \\ \\

    \theutterance \stepcounter{utterance}  
    & & \multicolumn{4}{p{0.6\linewidth}}{
        \cellcolor[rgb]{0.9,0.9,0.9}{
            \makecell[{{p{\linewidth}}}]{
                \texttt{\tiny{[P1$\langle$GM]}}
                \texttt{Mexxejt 'I' għal (7,1) b'suċċess. Il{-}grilja aġġornata tiegħek tidher hekk:} \\
\texttt{```} \\
\texttt{ 1234567} \\
\texttt{╔═╤═══╤═╗} \\
\texttt{║W│◌◌◌│I║ 1} \\
\texttt{╟─┼──┬┘◌║ 2} \\
\texttt{║◌│◌◌│◌◌║ 3} \\
\texttt{╟─┴──┘H◌║ 4} \\
\texttt{║◌◌◌◌◌◌◌║ 5} \\
\texttt{║C◌◌◌◌◌◌║ 6} \\
\texttt{║◌◌◌◌T◌◌║ 7} \\
\texttt{╚═══════╝} \\
\texttt{```} \\
\texttt{Bħalissa qed tilagħbu r{-}round 4 minn massimu ta' 20.} \\
\texttt{Intom flimkien akkumulajtu 2 minn 12 penali.} \\
\texttt{Il{-}plejer l{-}ieħor mexxa oġġett fuq il{-}grilja tiegħu.} \\
\\ 
\\ 
\texttt{X'inhu l{-}kmand li jmiss tiegħek?} \\
            }
        }
    }
    & \\ \\

    \theutterance \stepcounter{utterance}  
    & \multicolumn{4}{p{0.6\linewidth}}{
        \cellcolor[rgb]{0.9,0.9,0.9}{
            \makecell[{{p{\linewidth}}}]{
                \texttt{\tiny{[P1$\rangle$GM]}}
                \texttt{ĊAQLEQ: T, (3, 1)} \\
            }
        }
    }
    & & \\ \\

    \theutterance \stepcounter{utterance}  
    & & & \multicolumn{2}{p{0.3\linewidth}}{
        \cellcolor[rgb]{0.9,0.9,0.9}{
            \makecell[{{p{\linewidth}}}]{
                \texttt{\tiny{[GM$|$GM]}}
                \texttt{ĊAQLEQ: T, (3, 1)} \\
            }
        }
    }
    & & \\ \\

    \theutterance \stepcounter{utterance}  
    & & & \multicolumn{2}{p{0.3\linewidth}}{
        \cellcolor[rgb]{0.9,0.9,0.9}{
            \makecell[{{p{\linewidth}}}]{
                \texttt{\tiny{[GM$|$GM]}}
                \texttt{T moved to (3, 1)} \\
            }
        }
    }
    & & \\ \\

    \theutterance \stepcounter{utterance}  
    & & \multicolumn{4}{p{0.6\linewidth}}{
        \cellcolor[rgb]{0.9,0.9,0.9}{
            \makecell[{{p{\linewidth}}}]{
                \texttt{\tiny{[GM$\rangle$P2]}}
                \texttt{Mexxejt 'T' għal (4,1) b'suċċess. Il{-}grilja aġġornata tiegħek tidher hekk:} \\
\texttt{```} \\
\texttt{ 1234567} \\
\texttt{╔═╤═══╤═╗} \\
\texttt{║W│IT◌│H║ 1} \\
\texttt{╟─┼──┬┘◌║ 2} \\
\texttt{║◌│◌◌│◌C║ 3} \\
\texttt{╟─┴──┘◌◌║ 4} \\
\texttt{║◌◌◌◌◌◌◌║ 5} \\
\texttt{║◌◌◌◌◌◌◌║ 6} \\
\texttt{║◌◌◌◌◌◌◌║ 7} \\
\texttt{╚═══════╝} \\
\texttt{```} \\
\texttt{Bħalissa qed tilagħbu r{-}round 5 minn massimu ta' 20.} \\
\texttt{Intom flimkien akkumulajtu 2 minn 12 penali.} \\
\texttt{Il{-}plejer l{-}ieħor mexxa oġġett fuq il{-}grilja tiegħu.} \\
\\ 
\\ 
\texttt{X'inhu l{-}kmand li jmiss tiegħek?} \\
            }
        }
    }
    & \\ \\

    \theutterance \stepcounter{utterance}  
    & & & \multicolumn{4}{p{0.6\linewidth}}{
        \cellcolor[rgb]{0.9,0.9,0.9}{
            \makecell[{{p{\linewidth}}}]{
                \texttt{\tiny{[GM$\langle$P2]}}
                \texttt{ĊAQLEQ: C, (5, 1)} \\
            }
        }
    }
     \\ \\

    \theutterance \stepcounter{utterance}  
    & & & \multicolumn{2}{p{0.3\linewidth}}{
        \cellcolor[rgb]{0.9,0.9,0.9}{
            \makecell[{{p{\linewidth}}}]{
                \texttt{\tiny{[GM$|$GM]}}
                \texttt{ĊAQLEQ: C, (5, 1)} \\
            }
        }
    }
    & & \\ \\

    \theutterance \stepcounter{utterance}  
    & & & \multicolumn{2}{p{0.3\linewidth}}{
        \cellcolor[rgb]{0.9,0.9,0.9}{
            \makecell[{{p{\linewidth}}}]{
                \texttt{\tiny{[GM$|$GM]}}
                \texttt{C moved to (5, 1)} \\
            }
        }
    }
    & & \\ \\

    \theutterance \stepcounter{utterance}  
    & & \multicolumn{4}{p{0.6\linewidth}}{
        \cellcolor[rgb]{0.9,0.9,0.9}{
            \makecell[{{p{\linewidth}}}]{
                \texttt{\tiny{[P1$\langle$GM]}}
                \texttt{Mexxejt 'T' għal (3,1) b'suċċess. Il{-}grilja aġġornata tiegħek tidher hekk:} \\
\texttt{```} \\
\texttt{ 1234567} \\
\texttt{╔═╤═══╤═╗} \\
\texttt{║W│T◌◌│I║ 1} \\
\texttt{╟─┼──┬┘◌║ 2} \\
\texttt{║◌│◌◌│◌◌║ 3} \\
\texttt{╟─┴──┘H◌║ 4} \\
\texttt{║◌◌◌◌◌◌◌║ 5} \\
\texttt{║C◌◌◌◌◌◌║ 6} \\
\texttt{║◌◌◌◌◌◌◌║ 7} \\
\texttt{╚═══════╝} \\
\texttt{```} \\
\texttt{Bħalissa qed tilagħbu r{-}round 5 minn massimu ta' 20.} \\
\texttt{Intom flimkien akkumulajtu 2 minn 12 penali.} \\
\texttt{Il{-}plejer l{-}ieħor mexxa oġġett fuq il{-}grilja tiegħu.} \\
\\ 
\\ 
\texttt{X'inhu l{-}kmand li jmiss tiegħek?} \\
            }
        }
    }
    & \\ \\

    \theutterance \stepcounter{utterance}  
    & \multicolumn{4}{p{0.6\linewidth}}{
        \cellcolor[rgb]{0.9,0.9,0.9}{
            \makecell[{{p{\linewidth}}}]{
                \texttt{\tiny{[P1$\rangle$GM]}}
                \texttt{ĊAQLEQ: C, (4, 1)} \\
            }
        }
    }
    & & \\ \\

    \theutterance \stepcounter{utterance}  
    & & & \multicolumn{2}{p{0.3\linewidth}}{
        \cellcolor[rgb]{0.9,0.9,0.9}{
            \makecell[{{p{\linewidth}}}]{
                \texttt{\tiny{[GM$|$GM]}}
                \texttt{ĊAQLEQ: C, (4, 1)} \\
            }
        }
    }
    & & \\ \\

    \theutterance \stepcounter{utterance}  
    & & & \multicolumn{2}{p{0.3\linewidth}}{
        \cellcolor[rgb]{0.9,0.9,0.9}{
            \makecell[{{p{\linewidth}}}]{
                \texttt{\tiny{[GM$|$GM]}}
                \texttt{C moved to (4, 1)} \\
            }
        }
    }
    & & \\ \\

    \theutterance \stepcounter{utterance}  
    & & \multicolumn{4}{p{0.6\linewidth}}{
        \cellcolor[rgb]{0.9,0.9,0.9}{
            \makecell[{{p{\linewidth}}}]{
                \texttt{\tiny{[GM$\rangle$P2]}}
                \texttt{Mexxejt 'C' għal (5,1) b'suċċess. Il{-}grilja aġġornata tiegħek tidher hekk:} \\
\texttt{```} \\
\texttt{ 1234567} \\
\texttt{╔═╤═══╤═╗} \\
\texttt{║W│ITC│H║ 1} \\
\texttt{╟─┼──┬┘◌║ 2} \\
\texttt{║◌│◌◌│◌◌║ 3} \\
\texttt{╟─┴──┘◌◌║ 4} \\
\texttt{║◌◌◌◌◌◌◌║ 5} \\
\texttt{║◌◌◌◌◌◌◌║ 6} \\
\texttt{║◌◌◌◌◌◌◌║ 7} \\
\texttt{╚═══════╝} \\
\texttt{```} \\
\texttt{Bħalissa qed tilagħbu r{-}round 6 minn massimu ta' 20.} \\
\texttt{Intom flimkien akkumulajtu 2 minn 12 penali.} \\
\texttt{Il{-}plejer l{-}ieħor mexxa oġġett fuq il{-}grilja tiegħu.} \\
\\ 
\\ 
\texttt{X'inhu l{-}kmand li jmiss tiegħek?} \\
            }
        }
    }
    & \\ \\

    \theutterance \stepcounter{utterance}  
    & & & \multicolumn{4}{p{0.6\linewidth}}{
        \cellcolor[rgb]{0.9,0.9,0.9}{
            \makecell[{{p{\linewidth}}}]{
                \texttt{\tiny{[GM$\langle$P2]}}
                \texttt{ĊAQLEQ: H, (7, 1)} \\
            }
        }
    }
     \\ \\

    \theutterance \stepcounter{utterance}  
    & & & \multicolumn{2}{p{0.3\linewidth}}{
        \cellcolor[rgb]{0.9,0.9,0.9}{
            \makecell[{{p{\linewidth}}}]{
                \texttt{\tiny{[GM$|$GM]}}
                \texttt{ĊAQLEQ: H, (7, 1)} \\
            }
        }
    }
    & & \\ \\

    \theutterance \stepcounter{utterance}  
    & & & \multicolumn{2}{p{0.3\linewidth}}{
        \cellcolor[rgb]{0.9,0.9,0.9}{
            \makecell[{{p{\linewidth}}}]{
                \texttt{\tiny{[GM$|$GM]}}
                \texttt{Penali: (7,1) mhijiex vojta, iżda fiha 'H'.} \\
\\ 
\texttt{Intom flimkien akkumulajtu 3 minn 12 penali.} \\
\texttt{Jekk jogħġbok erġa' pprova!} \\
            }
        }
    }
    & & \\ \\

    \theutterance \stepcounter{utterance}  
    & & \multicolumn{4}{p{0.6\linewidth}}{
        \cellcolor[rgb]{0.9,0.9,0.9}{
            \makecell[{{p{\linewidth}}}]{
                \texttt{\tiny{[GM$\rangle$P2]}}
                \texttt{Penali: (7,1) mhijiex vojta, iżda fiha 'H'.} \\
\\ 
\texttt{Intom flimkien akkumulajtu 3 minn 12 penali.} \\
\texttt{Jekk jogħġbok erġa' pprova!} \\
            }
        }
    }
    & \\ \\

    \theutterance \stepcounter{utterance}  
    & & & \multicolumn{4}{p{0.6\linewidth}}{
        \cellcolor[rgb]{0.9,0.9,0.9}{
            \makecell[{{p{\linewidth}}}]{
                \texttt{\tiny{[GM$\langle$P2]}}
                \texttt{ĊAQLEQ: H, (6, 1)} \\
            }
        }
    }
     \\ \\

    \theutterance \stepcounter{utterance}  
    & & & \multicolumn{2}{p{0.3\linewidth}}{
        \cellcolor[rgb]{0.9,0.9,0.9}{
            \makecell[{{p{\linewidth}}}]{
                \texttt{\tiny{[GM$|$GM]}}
                \texttt{ĊAQLEQ: H, (6, 1)} \\
            }
        }
    }
    & & \\ \\

    \theutterance \stepcounter{utterance}  
    & & & \multicolumn{2}{p{0.3\linewidth}}{
        \cellcolor[rgb]{0.9,0.9,0.9}{
            \makecell[{{p{\linewidth}}}]{
                \texttt{\tiny{[GM$|$GM]}}
                \texttt{Penali: (6,1) mhijiex vojta, iżda fiha '│'.} \\
\\ 
\texttt{Intom flimkien akkumulajtu 4 minn 12 penali.} \\
\texttt{Jekk jogħġbok erġa' pprova!} \\
            }
        }
    }
    & & \\ \\

    \theutterance \stepcounter{utterance}  
    & & \multicolumn{4}{p{0.6\linewidth}}{
        \cellcolor[rgb]{0.9,0.9,0.9}{
            \makecell[{{p{\linewidth}}}]{
                \texttt{\tiny{[GM$\rangle$P2]}}
                \texttt{Penali: (6,1) mhijiex vojta, iżda fiha '│'.} \\
\\ 
\texttt{Intom flimkien akkumulajtu 4 minn 12 penali.} \\
\texttt{Jekk jogħġbok erġa' pprova!} \\
            }
        }
    }
    & \\ \\

    \theutterance \stepcounter{utterance}  
    & & & \multicolumn{4}{p{0.6\linewidth}}{
        \cellcolor[rgb]{0.9,0.9,0.9}{
            \makecell[{{p{\linewidth}}}]{
                \texttt{\tiny{[GM$\langle$P2]}}
                \texttt{GĦID: Fejn qegħdin l{-}oġġetti tiegħek issa? Għandi bżonn issa nċaqlaq 'H' għal (7,1) għax għandi 'H' dejjem f'(7,1). Huwa 'H' tiegħek qed jibqa f'(7,4) jew huwa mexxi? Jekk 'H' tiegħek f'(7,1) awtomatikament, nista' nċaqlaq tiegħi 'H' għal (7,1) u nkompletu l{-}għan. Mela fejn qegħdin l{-}oġġetti tiegħek issa?} \\
            }
        }
    }
     \\ \\

    \theutterance \stepcounter{utterance}  
    & & & \multicolumn{2}{p{0.3\linewidth}}{
        \cellcolor[rgb]{0.9,0.9,0.9}{
            \makecell[{{p{\linewidth}}}]{
                \texttt{\tiny{[GM$|$GM]}}
                \texttt{GĦID: Fejn qegħdin l{-}oġġetti tiegħek issa? Għandi bżonn issa nċaqlaq 'H' għal (7,1) għax għandi 'H' dejjem f'(7,1). Huwa 'H' tiegħek qed jibqa f'(7,4) jew huwa mexxi? Jekk 'H' tiegħek f'(7,1) awtomatikament, nista' nċaqlaq tiegħi 'H' għal (7,1) u nkompletu l{-}għan. Mela fejn qegħdin l{-}oġġetti tiegħek issa?} \\
            }
        }
    }
    & & \\ \\

    \theutterance \stepcounter{utterance}  
    & & \multicolumn{4}{p{0.6\linewidth}}{
        \cellcolor[rgb]{0.9,0.9,0.9}{
            \makecell[{{p{\linewidth}}}]{
                \texttt{\tiny{[P1$\langle$GM]}}
                \texttt{Mexxejt 'C' għal (4,1) b'suċċess. Il{-}grilja aġġornata tiegħek tidher hekk:} \\
\texttt{```} \\
\texttt{ 1234567} \\
\texttt{╔═╤═══╤═╗} \\
\texttt{║W│TC◌│I║ 1} \\
\texttt{╟─┼──┬┘◌║ 2} \\
\texttt{║◌│◌◌│◌◌║ 3} \\
\texttt{╟─┴──┘H◌║ 4} \\
\texttt{║◌◌◌◌◌◌◌║ 5} \\
\texttt{║◌◌◌◌◌◌◌║ 6} \\
\texttt{║◌◌◌◌◌◌◌║ 7} \\
\texttt{╚═══════╝} \\
\texttt{```} \\
\texttt{Bħalissa qed tilagħbu r{-}round 6 minn massimu ta' 20.} \\
\texttt{Intom flimkien akkumulajtu 4 minn 12 penali.} \\
\texttt{Il{-}plejer l{-}ieħor bagħat dan il{-}messaġġ:} \\
\\ 
\\ 
\texttt{"Fejn qegħdin l{-}oġġetti tiegħek issa? Għandi bżonn issa nċaqlaq 'H' għal (7,1) għax għandi 'H' dejjem f'(7,1). Huwa 'H' tiegħek qed jibqa f'(7,4) jew huwa mexxi? Jekk 'H' tiegħek f'(7,1) awtomatikament, nista' nċaqlaq tiegħi 'H' għal (7,1) u nkompletu l{-}għan. Mela fejn qegħdin l{-}oġġetti tiegħek issa?"} \\
\\ 
\\ 
\texttt{X'inhu l{-}kmand li jmiss tiegħek?} \\
            }
        }
    }
    & \\ \\

    \theutterance \stepcounter{utterance}  
    & \multicolumn{4}{p{0.6\linewidth}}{
        \cellcolor[rgb]{0.9,0.9,0.9}{
            \makecell[{{p{\linewidth}}}]{
                \texttt{\tiny{[P1$\rangle$GM]}}
                \texttt{GĦID: L{-}oġġetti tiegħi issa huma: W f'(1,1), I f'(7,1), T f'(3,1), C f'(4,1), H f'(7,4). Għax il{-}moviment għal (7,1) kien miċċaqqlaq mill{-}mur, rrid nigri l{-}Ħ f'(7,4) għal (5,1) sabiex nkompletu r{-}riga. Se nagħmel dan issa.} \\
            }
        }
    }
    & & \\ \\

    \theutterance \stepcounter{utterance}  
    & & & \multicolumn{2}{p{0.3\linewidth}}{
        \cellcolor[rgb]{0.9,0.9,0.9}{
            \makecell[{{p{\linewidth}}}]{
                \texttt{\tiny{[GM$|$GM]}}
                \texttt{GĦID: L{-}oġġetti tiegħi issa huma: W f'(1,1), I f'(7,1), T f'(3,1), C f'(4,1), H f'(7,4). Għax il{-}moviment għal (7,1) kien miċċaqqlaq mill{-}mur, rrid nigri l{-}Ħ f'(7,4) għal (5,1) sabiex nkompletu r{-}riga. Se nagħmel dan issa.} \\
            }
        }
    }
    & & \\ \\

    \theutterance \stepcounter{utterance}  
    & & \multicolumn{4}{p{0.6\linewidth}}{
        \cellcolor[rgb]{0.9,0.9,0.9}{
            \makecell[{{p{\linewidth}}}]{
                \texttt{\tiny{[GM$\rangle$P2]}}
                \texttt{Il{-}messaġġ tiegħek ġie mgħoddi lill{-}plejer l{-}ieħor.} \\
\\ 
\\ 
\texttt{Bħalissa qed tilagħbu r{-}round 7 minn massimu ta' 20.} \\
\texttt{Intom flimkien akkumulajtu 4 minn 12 penali.} \\
\texttt{Il{-}plejer l{-}ieħor bagħat dan il{-}messaġġ:} \\
\\ 
\\ 
\texttt{"L{-}oġġetti tiegħi issa huma: W f'(1,1), I f'(7,1), T f'(3,1), C f'(4,1), H f'(7,4). Għax il{-}moviment għal (7,1) kien miċċaqqlaq mill{-}mur, rrid nigri l{-}Ħ f'(7,4) għal (5,1) sabiex nkompletu r{-}riga. Se nagħmel dan issa."} \\
\\ 
\\ 
\texttt{X'inhu l{-}kmand li jmiss tiegħek?} \\
            }
        }
    }
    & \\ \\

    \theutterance \stepcounter{utterance}  
    & & & \multicolumn{4}{p{0.6\linewidth}}{
        \cellcolor[rgb]{0.9,0.9,0.9}{
            \makecell[{{p{\linewidth}}}]{
                \texttt{\tiny{[GM$\langle$P2]}}
                \texttt{GĦID: Mhux meħtieġ li tċaqlaq l{-}Ħ tiegħek għal (5,1) għax C tiegħek qed jinsab f'(4,1) u l{-}oġġetti ma jistgħux jibqgħu fl{-}istess pożizzjoni. Inħoss li kien hemm malfhem: l{-}għan huwa li l{-}oġġetti identiċi (W, I, T, C, H) jispiċċaw fl{-}istess koordinati fuq l{-}grilji tiegħna tnejn. Dak li naqbadna qabel (Riga 1: W(1,1), I(2,1), T(3,1), C(4,1), H(5,1)) huwa l{-}għan komuni li inti spiċċajt parti minnu (W, T, C f'pożizzjonijiet korretti). Għalhekk, għandna biss bżonn li inti tċaqlaq I tiegħek għal (2,1) u H tiegħek għal (5,1) (meta tkun vojta). Jien se nistenniex li tkompletu dak li qed nagħmel. Huwa dan l{-}akkord tiegħek?} \\
            }
        }
    }
     \\ \\

    \theutterance \stepcounter{utterance}  
    & & & \multicolumn{2}{p{0.3\linewidth}}{
        \cellcolor[rgb]{0.9,0.9,0.9}{
            \makecell[{{p{\linewidth}}}]{
                \texttt{\tiny{[GM$|$GM]}}
                \texttt{GĦID: Mhux meħtieġ li tċaqlaq l{-}Ħ tiegħek għal (5,1) għax C tiegħek qed jinsab f'(4,1) u l{-}oġġetti ma jistgħux jibqgħu fl{-}istess pożizzjoni. Inħoss li kien hemm malfhem: l{-}għan huwa li l{-}oġġetti identiċi (W, I, T, C, H) jispiċċaw fl{-}istess koordinati fuq l{-}grilji tiegħna tnejn. Dak li naqbadna qabel (Riga 1: W(1,1), I(2,1), T(3,1), C(4,1), H(5,1)) huwa l{-}għan komuni li inti spiċċajt parti minnu (W, T, C f'pożizzjonijiet korretti). Għalhekk, għandna biss bżonn li inti tċaqlaq I tiegħek għal (2,1) u H tiegħek għal (5,1) (meta tkun vojta). Jien se nistenniex li tkompletu dak li qed nagħmel. Huwa dan l{-}akkord tiegħek?} \\
            }
        }
    }
    & & \\ \\

    \theutterance \stepcounter{utterance}  
    & & \multicolumn{4}{p{0.6\linewidth}}{
        \cellcolor[rgb]{0.9,0.9,0.9}{
            \makecell[{{p{\linewidth}}}]{
                \texttt{\tiny{[P1$\langle$GM]}}
                \texttt{Il{-}messaġġ tiegħek ġie mgħoddi lill{-}plejer l{-}ieħor.} \\
\\ 
\\ 
\texttt{Bħalissa qed tilagħbu r{-}round 7 minn massimu ta' 20.} \\
\texttt{Intom flimkien akkumulajtu 4 minn 12 penali.} \\
\texttt{Il{-}plejer l{-}ieħor bagħat dan il{-}messaġġ:} \\
\\ 
\\ 
\texttt{"Mhux meħtieġ li tċaqlaq l{-}Ħ tiegħek għal (5,1) għax C tiegħek qed jinsab f'(4,1) u l{-}oġġetti ma jistgħux jibqgħu fl{-}istess pożizzjoni. Inħoss li kien hemm malfhem: l{-}għan huwa li l{-}oġġetti identiċi (W, I, T, C, H) jispiċċaw fl{-}istess koordinati fuq l{-}grilji tiegħna tnejn. Dak li naqbadna qabel (Riga 1: W(1,1), I(2,1), T(3,1), C(4,1), H(5,1)) huwa l{-}għan komuni li inti spiċċajt parti minnu (W, T, C f'pożizzjonijiet korretti). Għalhekk, għandna biss bżonn li inti tċaqlaq I tiegħek għal (2,1) u H tiegħek għal (5,1) (meta tkun vojta). Jien se nistenniex li tkompletu dak li qed nagħmel. Huwa dan l{-}akkord tiegħek?"} \\
\\ 
\\ 
\texttt{X'inhu l{-}kmand li jmiss tiegħek?} \\
            }
        }
    }
    & \\ \\

    \theutterance \stepcounter{utterance}  
    & \multicolumn{4}{p{0.6\linewidth}}{
        \cellcolor[rgb]{0.9,0.9,0.9}{
            \makecell[{{p{\linewidth}}}]{
                \texttt{\tiny{[P1$\rangle$GM]}}
                \texttt{GĦID: Tajjeb, għandek raġuni. L{-}għan huwa li l{-}oġġetti identiċi jispiċċaw fl{-}istess koordinati. Intimappost għax kien hemm 'I' f'(7,1) u mhux f'(2,1). Se nigri 'H' tiegħi għal (5,1) u 'I' tiegħi għal (2,1) meta tkun vojta. Int qed tistenni, mhux hekk?} \\
            }
        }
    }
    & & \\ \\

    \theutterance \stepcounter{utterance}  
    & & & \multicolumn{2}{p{0.3\linewidth}}{
        \cellcolor[rgb]{0.9,0.9,0.9}{
            \makecell[{{p{\linewidth}}}]{
                \texttt{\tiny{[GM$|$GM]}}
                \texttt{GĦID: Tajjeb, għandek raġuni. L{-}għan huwa li l{-}oġġetti identiċi jispiċċaw fl{-}istess koordinati. Intimappost għax kien hemm 'I' f'(7,1) u mhux f'(2,1). Se nigri 'H' tiegħi għal (5,1) u 'I' tiegħi għal (2,1) meta tkun vojta. Int qed tistenni, mhux hekk?} \\
            }
        }
    }
    & & \\ \\

    \theutterance \stepcounter{utterance}  
    & & \multicolumn{4}{p{0.6\linewidth}}{
        \cellcolor[rgb]{0.9,0.9,0.9}{
            \makecell[{{p{\linewidth}}}]{
                \texttt{\tiny{[GM$\rangle$P2]}}
                \texttt{Il{-}messaġġ tiegħek ġie mgħoddi lill{-}plejer l{-}ieħor.} \\
\\ 
\\ 
\texttt{Bħalissa qed tilagħbu r{-}round 8 minn massimu ta' 20.} \\
\texttt{Intom flimkien akkumulajtu 4 minn 12 penali.} \\
\texttt{Il{-}plejer l{-}ieħor bagħat dan il{-}messaġġ:} \\
\\ 
\\ 
\texttt{"Tajjeb, għandek raġuni. L{-}għan huwa li l{-}oġġetti identiċi jispiċċaw fl{-}istess koordinati. Intimappost għax kien hemm 'I' f'(7,1) u mhux f'(2,1). Se nigri 'H' tiegħi għal (5,1) u 'I' tiegħi għal (2,1) meta tkun vojta. Int qed tistenni, mhux hekk?"} \\
\\ 
\\ 
\texttt{X'inhu l{-}kmand li jmiss tiegħek?} \\
            }
        }
    }
    & \\ \\

    \theutterance \stepcounter{utterance}  
    & & & \multicolumn{4}{p{0.6\linewidth}}{
        \cellcolor[rgb]{0.9,0.9,0.9}{
            \makecell[{{p{\linewidth}}}]{
                \texttt{\tiny{[GM$\langle$P2]}}
                \texttt{GĦID: Iva, tajjeb. Jien nistenni meta tkun (5,1) u (2,1) vojta għaliex għandi 'H' f'(7,1) u 'I' f'(4,6) li għandhom jibqgħu fejn huma. Jekk inti tista' tċaqlaq 'H' tiegħek għal (5,1) u 'I' tiegħek għal (2,1) bla ma tenniesa l{-}oġġetti tiegħi oħra, nkompletu l{-}għan. Nifhem li tista' tmexxi oġġett wieħed kull rawnd, għalhekk għandna bżonn rawndejn aktar. Se nkomunikaw meta tista' tagħmel il{-}movimenti.} \\
            }
        }
    }
     \\ \\

    \theutterance \stepcounter{utterance}  
    & & & \multicolumn{2}{p{0.3\linewidth}}{
        \cellcolor[rgb]{0.9,0.9,0.9}{
            \makecell[{{p{\linewidth}}}]{
                \texttt{\tiny{[GM$|$GM]}}
                \texttt{GĦID: Iva, tajjeb. Jien nistenni meta tkun (5,1) u (2,1) vojta għaliex għandi 'H' f'(7,1) u 'I' f'(4,6) li għandhom jibqgħu fejn huma. Jekk inti tista' tċaqlaq 'H' tiegħek għal (5,1) u 'I' tiegħek għal (2,1) bla ma tenniesa l{-}oġġetti tiegħi oħra, nkompletu l{-}għan. Nifhem li tista' tmexxi oġġett wieħed kull rawnd, għalhekk għandna bżonn rawndejn aktar. Se nkomunikaw meta tista' tagħmel il{-}movimenti.} \\
            }
        }
    }
    & & \\ \\

    \theutterance \stepcounter{utterance}  
    & & \multicolumn{4}{p{0.6\linewidth}}{
        \cellcolor[rgb]{0.9,0.9,0.9}{
            \makecell[{{p{\linewidth}}}]{
                \texttt{\tiny{[P1$\langle$GM]}}
                \texttt{Il{-}messaġġ tiegħek ġie mgħoddi lill{-}plejer l{-}ieħor.} \\
\\ 
\\ 
\texttt{Bħalissa qed tilagħbu r{-}round 8 minn massimu ta' 20.} \\
\texttt{Intom flimkien akkumulajtu 4 minn 12 penali.} \\
\texttt{Il{-}plejer l{-}ieħor bagħat dan il{-}messaġġ:} \\
\\ 
\\ 
\texttt{"Iva, tajjeb. Jien nistenni meta tkun (5,1) u (2,1) vojta għaliex għandi 'H' f'(7,1) u 'I' f'(4,6) li għandhom jibqgħu fejn huma. Jekk inti tista' tċaqlaq 'H' tiegħek għal (5,1) u 'I' tiegħek għal (2,1) bla ma tenniesa l{-}oġġetti tiegħi oħra, nkompletu l{-}għan. Nifhem li tista' tmexxi oġġett wieħed kull rawnd, għalhekk għandna bżonn rawndejn aktar. Se nkomunikaw meta tista' tagħmel il{-}movimenti."} \\
\\ 
\\ 
\texttt{X'inhu l{-}kmand li jmiss tiegħek?} \\
            }
        }
    }
    & \\ \\

    \theutterance \stepcounter{utterance}  
    & \multicolumn{4}{p{0.6\linewidth}}{
        \cellcolor[rgb]{0.9,0.9,0.9}{
            \makecell[{{p{\linewidth}}}]{
                \texttt{\tiny{[P1$\rangle$GM]}}
                \texttt{ĊAQLEQ: H, (5, 1)} \\
            }
        }
    }
    & & \\ \\

    \theutterance \stepcounter{utterance}  
    & & & \multicolumn{2}{p{0.3\linewidth}}{
        \cellcolor[rgb]{0.9,0.9,0.9}{
            \makecell[{{p{\linewidth}}}]{
                \texttt{\tiny{[GM$|$GM]}}
                \texttt{ĊAQLEQ: H, (5, 1)} \\
            }
        }
    }
    & & \\ \\

    \theutterance \stepcounter{utterance}  
    & & & \multicolumn{2}{p{0.3\linewidth}}{
        \cellcolor[rgb]{0.9,0.9,0.9}{
            \makecell[{{p{\linewidth}}}]{
                \texttt{\tiny{[GM$|$GM]}}
                \texttt{H moved to (5, 1)} \\
            }
        }
    }
    & & \\ \\

    \theutterance \stepcounter{utterance}  
    & & \multicolumn{4}{p{0.6\linewidth}}{
        \cellcolor[rgb]{0.9,0.9,0.9}{
            \makecell[{{p{\linewidth}}}]{
                \texttt{\tiny{[GM$\rangle$P2]}}
                \texttt{Il{-}messaġġ tiegħek ġie mgħoddi lill{-}plejer l{-}ieħor.} \\
\\ 
\\ 
\texttt{Bħalissa qed tilagħbu r{-}round 9 minn massimu ta' 20.} \\
\texttt{Intom flimkien akkumulajtu 4 minn 12 penali.} \\
\texttt{Il{-}plejer l{-}ieħor mexxa oġġett fuq il{-}grilja tiegħu.} \\
\\ 
\\ 
\texttt{X'inhu l{-}kmand li jmiss tiegħek?} \\
            }
        }
    }
    & \\ \\

    \theutterance \stepcounter{utterance}  
    & & & \multicolumn{4}{p{0.6\linewidth}}{
        \cellcolor[rgb]{0.9,0.9,0.9}{
            \makecell[{{p{\linewidth}}}]{
                \texttt{\tiny{[GM$\langle$P2]}}
                \texttt{GĦID: Kif sewwa? Int mexxejt 'H' tiegħek għal (5,1) jew 'I' tiegħek għal (2,1)? Fejn qegħdin l{-}oġġetti tiegħek issa?} \\
            }
        }
    }
     \\ \\

    \theutterance \stepcounter{utterance}  
    & & & \multicolumn{2}{p{0.3\linewidth}}{
        \cellcolor[rgb]{0.9,0.9,0.9}{
            \makecell[{{p{\linewidth}}}]{
                \texttt{\tiny{[GM$|$GM]}}
                \texttt{GĦID: Kif sewwa? Int mexxejt 'H' tiegħek għal (5,1) jew 'I' tiegħek għal (2,1)? Fejn qegħdin l{-}oġġetti tiegħek issa?} \\
            }
        }
    }
    & & \\ \\

    \theutterance \stepcounter{utterance}  
    & & \multicolumn{4}{p{0.6\linewidth}}{
        \cellcolor[rgb]{0.9,0.9,0.9}{
            \makecell[{{p{\linewidth}}}]{
                \texttt{\tiny{[P1$\langle$GM]}}
                \texttt{Mexxejt 'H' għal (5,1) b'suċċess. Il{-}grilja aġġornata tiegħek tidher hekk:} \\
\texttt{```} \\
\texttt{ 1234567} \\
\texttt{╔═╤═══╤═╗} \\
\texttt{║W│TCH│I║ 1} \\
\texttt{╟─┼──┬┘◌║ 2} \\
\texttt{║◌│◌◌│◌◌║ 3} \\
\texttt{╟─┴──┘◌◌║ 4} \\
\texttt{║◌◌◌◌◌◌◌║ 5} \\
\texttt{║◌◌◌◌◌◌◌║ 6} \\
\texttt{║◌◌◌◌◌◌◌║ 7} \\
\texttt{╚═══════╝} \\
\texttt{```} \\
\texttt{Bħalissa qed tilagħbu r{-}round 9 minn massimu ta' 20.} \\
\texttt{Intom flimkien akkumulajtu 4 minn 12 penali.} \\
\texttt{Il{-}plejer l{-}ieħor bagħat dan il{-}messaġġ:} \\
\\ 
\\ 
\texttt{"Kif sewwa? Int mexxejt 'H' tiegħek għal (5,1) jew 'I' tiegħek għal (2,1)? Fejn qegħdin l{-}oġġetti tiegħek issa?"} \\
\\ 
\\ 
\texttt{X'inhu l{-}kmand li jmiss tiegħek?} \\
            }
        }
    }
    & \\ \\

    \theutterance \stepcounter{utterance}  
    & \multicolumn{4}{p{0.6\linewidth}}{
        \cellcolor[rgb]{0.9,0.9,0.9}{
            \makecell[{{p{\linewidth}}}]{
                \texttt{\tiny{[P1$\rangle$GM]}}
                \texttt{ĊAQLEQ: I, (2, 1)} \\
            }
        }
    }
    & & \\ \\

    \theutterance \stepcounter{utterance}  
    & & & \multicolumn{2}{p{0.3\linewidth}}{
        \cellcolor[rgb]{0.9,0.9,0.9}{
            \makecell[{{p{\linewidth}}}]{
                \texttt{\tiny{[GM$|$GM]}}
                \texttt{ĊAQLEQ: I, (2, 1)} \\
            }
        }
    }
    & & \\ \\

    \theutterance \stepcounter{utterance}  
    & & & \multicolumn{2}{p{0.3\linewidth}}{
        \cellcolor[rgb]{0.9,0.9,0.9}{
            \makecell[{{p{\linewidth}}}]{
                \texttt{\tiny{[GM$|$GM]}}
                \texttt{Penali: (2,1) mhijiex vojta, iżda fiha '│'.} \\
\\ 
\texttt{Intom flimkien akkumulajtu 5 minn 12 penali.} \\
\texttt{Jekk jogħġbok erġa' pprova!} \\
            }
        }
    }
    & & \\ \\

    \theutterance \stepcounter{utterance}  
    & & \multicolumn{4}{p{0.6\linewidth}}{
        \cellcolor[rgb]{0.9,0.9,0.9}{
            \makecell[{{p{\linewidth}}}]{
                \texttt{\tiny{[P1$\langle$GM]}}
                \texttt{Penali: (2,1) mhijiex vojta, iżda fiha '│'.} \\
\\ 
\texttt{Intom flimkien akkumulajtu 5 minn 12 penali.} \\
\texttt{Jekk jogħġbok erġa' pprova!} \\
            }
        }
    }
    & \\ \\

    \theutterance \stepcounter{utterance}  
    & \multicolumn{4}{p{0.6\linewidth}}{
        \cellcolor[rgb]{0.9,0.9,0.9}{
            \makecell[{{p{\linewidth}}}]{
                \texttt{\tiny{[P1$\rangle$GM]}}
                \texttt{GĦID: L{-}istat attwali: W(1,1), I(7,1), T(3,1), C(4,1), H(5,1). Qed nistenn li (2,1) tispiċċa vojta. Int wasalt għall{-}konfigurazzjoni ta' għan?} \\
            }
        }
    }
    & & \\ \\

    \theutterance \stepcounter{utterance}  
    & & & \multicolumn{2}{p{0.3\linewidth}}{
        \cellcolor[rgb]{0.9,0.9,0.9}{
            \makecell[{{p{\linewidth}}}]{
                \texttt{\tiny{[GM$|$GM]}}
                \texttt{GĦID: L{-}istat attwali: W(1,1), I(7,1), T(3,1), C(4,1), H(5,1). Qed nistenn li (2,1) tispiċċa vojta. Int wasalt għall{-}konfigurazzjoni ta' għan?} \\
            }
        }
    }
    & & \\ \\

    \theutterance \stepcounter{utterance}  
    & & \multicolumn{4}{p{0.6\linewidth}}{
        \cellcolor[rgb]{0.9,0.9,0.9}{
            \makecell[{{p{\linewidth}}}]{
                \texttt{\tiny{[GM$\rangle$P2]}}
                \texttt{Il{-}messaġġ tiegħek ġie mgħoddi lill{-}plejer l{-}ieħor.} \\
\\ 
\\ 
\texttt{Bħalissa qed tilagħbu r{-}round 10 minn massimu ta' 20.} \\
\texttt{Intom flimkien akkumulajtu 5 minn 12 penali.} \\
\texttt{Il{-}plejer l{-}ieħor bagħat dan il{-}messaġġ:} \\
\\ 
\\ 
\texttt{"L{-}istat attwali: W(1,1), I(7,1), T(3,1), C(4,1), H(5,1). Qed nistenn li (2,1) tispiċċa vojta. Int wasalt għall{-}konfigurazzjoni ta' għan?"} \\
\\ 
\\ 
\texttt{X'inhu l{-}kmand li jmiss tiegħek?} \\
            }
        }
    }
    & \\ \\

    \theutterance \stepcounter{utterance}  
    & & & \multicolumn{4}{p{0.6\linewidth}}{
        \cellcolor[rgb]{0.9,0.9,0.9}{
            \makecell[{{p{\linewidth}}}]{
                \texttt{\tiny{[GM$\langle$P2]}}
                \texttt{GĦID: Le, għajda mhux wasalt għall{-}konfigurazzjoni ta' għan. L{-}oġġetti tiegħi huma: W(1,1), I(3,1), T(4,1), C(5,1), H(7,1). Għandi bżonn li inti tċaqlaq 'I' tiegħek minn (7,1) għal (2,1) sabiex nispiċċaw. Meta tagħmel dan, nkomplu. Fejn qegħdin l{-}oġġetti tiegħek issa? Qed tistenni li (2,1) tikber vojta?} \\
            }
        }
    }
     \\ \\

    \theutterance \stepcounter{utterance}  
    & & & \multicolumn{2}{p{0.3\linewidth}}{
        \cellcolor[rgb]{0.9,0.9,0.9}{
            \makecell[{{p{\linewidth}}}]{
                \texttt{\tiny{[GM$|$GM]}}
                \texttt{GĦID: Le, għajda mhux wasalt għall{-}konfigurazzjoni ta' għan. L{-}oġġetti tiegħi huma: W(1,1), I(3,1), T(4,1), C(5,1), H(7,1). Għandi bżonn li inti tċaqlaq 'I' tiegħek minn (7,1) għal (2,1) sabiex nispiċċaw. Meta tagħmel dan, nkomplu. Fejn qegħdin l{-}oġġetti tiegħek issa? Qed tistenni li (2,1) tikber vojta?} \\
            }
        }
    }
    & & \\ \\

    \theutterance \stepcounter{utterance}  
    & & \multicolumn{4}{p{0.6\linewidth}}{
        \cellcolor[rgb]{0.9,0.9,0.9}{
            \makecell[{{p{\linewidth}}}]{
                \texttt{\tiny{[P1$\langle$GM]}}
                \texttt{Il{-}messaġġ tiegħek ġie mgħoddi lill{-}plejer l{-}ieħor.} \\
\\ 
\\ 
\texttt{Bħalissa qed tilagħbu r{-}round 10 minn massimu ta' 20.} \\
\texttt{Intom flimkien akkumulajtu 5 minn 12 penali.} \\
\texttt{Il{-}plejer l{-}ieħor bagħat dan il{-}messaġġ:} \\
\\ 
\\ 
\texttt{"Le, għajda mhux wasalt għall{-}konfigurazzjoni ta' għan. L{-}oġġetti tiegħi huma: W(1,1), I(3,1), T(4,1), C(5,1), H(7,1). Għandi bżonn li inti tċaqlaq 'I' tiegħek minn (7,1) għal (2,1) sabiex nispiċċaw. Meta tagħmel dan, nkomplu. Fejn qegħdin l{-}oġġetti tiegħek issa? Qed tistenni li (2,1) tikber vojta?"} \\
\\ 
\\ 
\texttt{X'inhu l{-}kmand li jmiss tiegħek?} \\
            }
        }
    }
    & \\ \\

    \theutterance \stepcounter{utterance}  
    & \multicolumn{4}{p{0.6\linewidth}}{
        \cellcolor[rgb]{0.9,0.9,0.9}{
            \makecell[{{p{\linewidth}}}]{
                \texttt{\tiny{[P1$\rangle$GM]}}
                \texttt{GĦID: L{-}oġġetti tiegħi: W(1,1), I(7,1), T(3,1), C(4,1), H(5,1). Għandi bżonn inħalli (7,1) vojta qabel ma nista' nċaqlaq I għal (2,1). X'kmand qed tagħmel int? Tista' tċaqlaq I tiegħek minn (4,6) għal (3,1)?} \\
            }
        }
    }
    & & \\ \\

    \theutterance \stepcounter{utterance}  
    & & & \multicolumn{2}{p{0.3\linewidth}}{
        \cellcolor[rgb]{0.9,0.9,0.9}{
            \makecell[{{p{\linewidth}}}]{
                \texttt{\tiny{[GM$|$GM]}}
                \texttt{GĦID: L{-}oġġetti tiegħi: W(1,1), I(7,1), T(3,1), C(4,1), H(5,1). Għandi bżonn inħalli (7,1) vojta qabel ma nista' nċaqlaq I għal (2,1). X'kmand qed tagħmel int? Tista' tċaqlaq I tiegħek minn (4,6) għal (3,1)?} \\
            }
        }
    }
    & & \\ \\

    \theutterance \stepcounter{utterance}  
    & & \multicolumn{4}{p{0.6\linewidth}}{
        \cellcolor[rgb]{0.9,0.9,0.9}{
            \makecell[{{p{\linewidth}}}]{
                \texttt{\tiny{[GM$\rangle$P2]}}
                \texttt{Il{-}messaġġ tiegħek ġie mgħoddi lill{-}plejer l{-}ieħor.} \\
\\ 
\\ 
\texttt{Bħalissa qed tilagħbu r{-}round 11 minn massimu ta' 20.} \\
\texttt{Intom flimkien akkumulajtu 5 minn 12 penali.} \\
\texttt{Il{-}plejer l{-}ieħor bagħat dan il{-}messaġġ:} \\
\\ 
\\ 
\texttt{"L{-}oġġetti tiegħi: W(1,1), I(7,1), T(3,1), C(4,1), H(5,1). Għandi bżonn inħalli (7,1) vojta qabel ma nista' nċaqlaq I għal (2,1). X'kmand qed tagħmel int? Tista' tċaqlaq I tiegħek minn (4,6) għal (3,1)?"} \\
\\ 
\\ 
\texttt{X'inhu l{-}kmand li jmiss tiegħek?} \\
            }
        }
    }
    & \\ \\

    \theutterance \stepcounter{utterance}  
    & & & \multicolumn{4}{p{0.6\linewidth}}{
        \cellcolor[rgb]{0.9,0.9,0.9}{
            \makecell[{{p{\linewidth}}}]{
                \texttt{\tiny{[GM$\langle$P2]}}
                \texttt{GĦID: Le, m'għandix 'I' f'(4,6) aktar. Mexxejt 'I' għal (3,1) f'rawnd 3. L{-}oġġetti tiegħi issa huma: W(1,1), I(3,1), T(4,1), C(5,1), H(7,1). Għandi bżonn li inti tċaqlaq 'I' tiegħek minn (7,1) għal (2,1) u 'H' tiegħek qed jinsab f'(5,1) (waqgħa dik mexxejt). Meta tagħmel dan, il{-}konfigurazzjoni tiegħi u tiegħek se tkun identika: W(1,1), I(2,1), T(3,1), C(4,1), H(5,1). Jekk inti mexxejt 'I' tiegħek għal (2,1) issa, nista' nibda nċaqlaq 'H' tiegħi għal (6,1) jew (5,1) meta tkun vojta, imma għandi bżonn li l{-}konfigurazzjoni tkun identika qabel. Fejn qegħdin l{-}oġġetti tiegħek ISSA? Hu 'I' tiegħek anqas f'(7,1)?} \\
            }
        }
    }
     \\ \\

    \theutterance \stepcounter{utterance}  
    & & & \multicolumn{2}{p{0.3\linewidth}}{
        \cellcolor[rgb]{0.9,0.9,0.9}{
            \makecell[{{p{\linewidth}}}]{
                \texttt{\tiny{[GM$|$GM]}}
                \texttt{GĦID: Le, m'għandix 'I' f'(4,6) aktar. Mexxejt 'I' għal (3,1) f'rawnd 3. L{-}oġġetti tiegħi issa huma: W(1,1), I(3,1), T(4,1), C(5,1), H(7,1). Għandi bżonn li inti tċaqlaq 'I' tiegħek minn (7,1) għal (2,1) u 'H' tiegħek qed jinsab f'(5,1) (waqgħa dik mexxejt). Meta tagħmel dan, il{-}konfigurazzjoni tiegħi u tiegħek se tkun identika: W(1,1), I(2,1), T(3,1), C(4,1), H(5,1). Jekk inti mexxejt 'I' tiegħek għal (2,1) issa, nista' nibda nċaqlaq 'H' tiegħi għal (6,1) jew (5,1) meta tkun vojta, imma għandi bżonn li l{-}konfigurazzjoni tkun identika qabel. Fejn qegħdin l{-}oġġetti tiegħek ISSA? Hu 'I' tiegħek anqas f'(7,1)?} \\
            }
        }
    }
    & & \\ \\

    \theutterance \stepcounter{utterance}  
    & & \multicolumn{4}{p{0.6\linewidth}}{
        \cellcolor[rgb]{0.9,0.9,0.9}{
            \makecell[{{p{\linewidth}}}]{
                \texttt{\tiny{[P1$\langle$GM]}}
                \texttt{Il{-}messaġġ tiegħek ġie mgħoddi lill{-}plejer l{-}ieħor.} \\
\\ 
\\ 
\texttt{Bħalissa qed tilagħbu r{-}round 11 minn massimu ta' 20.} \\
\texttt{Intom flimkien akkumulajtu 5 minn 12 penali.} \\
\texttt{Il{-}plejer l{-}ieħor bagħat dan il{-}messaġġ:} \\
\\ 
\\ 
\texttt{"Le, m'għandix 'I' f'(4,6) aktar. Mexxejt 'I' għal (3,1) f'rawnd 3. L{-}oġġetti tiegħi issa huma: W(1,1), I(3,1), T(4,1), C(5,1), H(7,1). Għandi bżonn li inti tċaqlaq 'I' tiegħek minn (7,1) għal (2,1) u 'H' tiegħek qed jinsab f'(5,1) (waqgħa dik mexxejt). Meta tagħmel dan, il{-}konfigurazzjoni tiegħi u tiegħek se tkun identika: W(1,1), I(2,1), T(3,1), C(4,1), H(5,1). Jekk inti mexxejt 'I' tiegħek għal (2,1) issa, nista' nibda nċaqlaq 'H' tiegħi għal (6,1) jew (5,1) meta tkun vojta, imma għandi bżonn li l{-}konfigurazzjoni tkun identika qabel. Fejn qegħdin l{-}oġġetti tiegħek ISSA? Hu 'I' tiegħek anqas f'(7,1)?"} \\
\\ 
\\ 
\texttt{X'inhu l{-}kmand li jmiss tiegħek?} \\
            }
        }
    }
    & \\ \\

    \theutterance \stepcounter{utterance}  
    & \multicolumn{4}{p{0.6\linewidth}}{
        \cellcolor[rgb]{0.9,0.9,0.9}{
            \makecell[{{p{\linewidth}}}]{
                \texttt{\tiny{[P1$\rangle$GM]}}
                \texttt{ĊAQLEQ: I, (2, 1)} \\
            }
        }
    }
    & & \\ \\

    \theutterance \stepcounter{utterance}  
    & & & \multicolumn{2}{p{0.3\linewidth}}{
        \cellcolor[rgb]{0.9,0.9,0.9}{
            \makecell[{{p{\linewidth}}}]{
                \texttt{\tiny{[GM$|$GM]}}
                \texttt{ĊAQLEQ: I, (2, 1)} \\
            }
        }
    }
    & & \\ \\

    \theutterance \stepcounter{utterance}  
    & & & \multicolumn{2}{p{0.3\linewidth}}{
        \cellcolor[rgb]{0.9,0.9,0.9}{
            \makecell[{{p{\linewidth}}}]{
                \texttt{\tiny{[GM$|$GM]}}
                \texttt{Penali: (2,1) mhijiex vojta, iżda fiha '│'.} \\
\\ 
\texttt{Intom flimkien akkumulajtu 6 minn 12 penali.} \\
\texttt{Jekk jogħġbok erġa' pprova!} \\
            }
        }
    }
    & & \\ \\

    \theutterance \stepcounter{utterance}  
    & & \multicolumn{4}{p{0.6\linewidth}}{
        \cellcolor[rgb]{0.9,0.9,0.9}{
            \makecell[{{p{\linewidth}}}]{
                \texttt{\tiny{[P1$\langle$GM]}}
                \texttt{Penali: (2,1) mhijiex vojta, iżda fiha '│'.} \\
\\ 
\texttt{Intom flimkien akkumulajtu 6 minn 12 penali.} \\
\texttt{Jekk jogħġbok erġa' pprova!} \\
            }
        }
    }
    & \\ \\

    \theutterance \stepcounter{utterance}  
    & \multicolumn{4}{p{0.6\linewidth}}{
        \cellcolor[rgb]{0.9,0.9,0.9}{
            \makecell[{{p{\linewidth}}}]{
                \texttt{\tiny{[P1$\rangle$GM]}}
                \texttt{GĦID: Il{-}kelma (2,1) mhix vojta għax fih '│' (ħajt vertikali). Ma nistax nċaqlaq 'I' hemm. Il{-}konfigurazzjoni li qed nitlobu (I f'(2,1)) mhix possibbli għax (2,1) hija ċella mblokkata mill{-}ħajtar. Għandna bżonn nitlobu għan ieħor li jikseb pożizzjonijiet vojta għal kulħadd. Fejn huma l{-}oġġetti tiegħek ISSA?} \\
            }
        }
    }
    & & \\ \\

    \theutterance \stepcounter{utterance}  
    & & & \multicolumn{2}{p{0.3\linewidth}}{
        \cellcolor[rgb]{0.9,0.9,0.9}{
            \makecell[{{p{\linewidth}}}]{
                \texttt{\tiny{[GM$|$GM]}}
                \texttt{GĦID: Il{-}kelma (2,1) mhix vojta għax fih '│' (ħajt vertikali). Ma nistax nċaqlaq 'I' hemm. Il{-}konfigurazzjoni li qed nitlobu (I f'(2,1)) mhix possibbli għax (2,1) hija ċella mblokkata mill{-}ħajtar. Għandna bżonn nitlobu għan ieħor li jikseb pożizzjonijiet vojta għal kulħadd. Fejn huma l{-}oġġetti tiegħek ISSA?} \\
            }
        }
    }
    & & \\ \\

    \theutterance \stepcounter{utterance}  
    & & \multicolumn{4}{p{0.6\linewidth}}{
        \cellcolor[rgb]{0.9,0.9,0.9}{
            \makecell[{{p{\linewidth}}}]{
                \texttt{\tiny{[GM$\rangle$P2]}}
                \texttt{Il{-}messaġġ tiegħek ġie mgħoddi lill{-}plejer l{-}ieħor.} \\
\\ 
\\ 
\texttt{Bħalissa qed tilagħbu r{-}round 12 minn massimu ta' 20.} \\
\texttt{Intom flimkien akkumulajtu 6 minn 12 penali.} \\
\texttt{Il{-}plejer l{-}ieħor bagħat dan il{-}messaġġ:} \\
\\ 
\\ 
\texttt{"Il{-}kelma (2,1) mhix vojta għax fih '│' (ħajt vertikali). Ma nistax nċaqlaq 'I' hemm. Il{-}konfigurazzjoni li qed nitlobu (I f'(2,1)) mhix possibbli għax (2,1) hija ċella mblokkata mill{-}ħajtar. Għandna bżonn nitlobu għan ieħor li jikseb pożizzjonijiet vojta għal kulħadd. Fejn huma l{-}oġġetti tiegħek ISSA?"} \\
\\ 
\\ 
\texttt{X'inhu l{-}kmand li jmiss tiegħek?} \\
            }
        }
    }
    & \\ \\

    \theutterance \stepcounter{utterance}  
    & & & \multicolumn{4}{p{0.6\linewidth}}{
        \cellcolor[rgb]{0.9,0.9,0.9}{
            \makecell[{{p{\linewidth}}}]{
                \texttt{\tiny{[GM$\langle$P2]}}
                \texttt{GĦID: Aħna qegħdin ninsa l{-}ħajtar f'(2,1)! Għaxda, l{-}għan li qed nitlobu mhux possibbli. Il{-}kordinati vojta f'Riga 1 huma: (1,1), (3,1), (4,1), (5,1), (6,1). (2,1) u (7,1) huma mblokkati. Għandna 5 oġġetti u 5 kordinati vojta {-} perfett! Nitlobu għan ġdid: W(1,1), I(3,1), T(4,1), C(5,1), H(6,1). Dak huwa l{-}istess li għandi issa! Int wasalt għal W(1,1), T(3,1), C(4,1), H(5,1). Għandek bżonn tċaqlaq I għal (6,1) u T għal (3,1) {-} tista' tagħmel dan? Għidli fejn qegħdin l{-}oġġetti tiegħek ISSA u jekk tista' tagħmel dan.} \\
            }
        }
    }
     \\ \\

    \theutterance \stepcounter{utterance}  
    & & & \multicolumn{2}{p{0.3\linewidth}}{
        \cellcolor[rgb]{0.9,0.9,0.9}{
            \makecell[{{p{\linewidth}}}]{
                \texttt{\tiny{[GM$|$GM]}}
                \texttt{GĦID: Aħna qegħdin ninsa l{-}ħajtar f'(2,1)! Għaxda, l{-}għan li qed nitlobu mhux possibbli. Il{-}kordinati vojta f'Riga 1 huma: (1,1), (3,1), (4,1), (5,1), (6,1). (2,1) u (7,1) huma mblokkati. Għandna 5 oġġetti u 5 kordinati vojta {-} perfett! Nitlobu għan ġdid: W(1,1), I(3,1), T(4,1), C(5,1), H(6,1). Dak huwa l{-}istess li għandi issa! Int wasalt għal W(1,1), T(3,1), C(4,1), H(5,1). Għandek bżonn tċaqlaq I għal (6,1) u T għal (3,1) {-} tista' tagħmel dan? Għidli fejn qegħdin l{-}oġġetti tiegħek ISSA u jekk tista' tagħmel dan.} \\
            }
        }
    }
    & & \\ \\

    \theutterance \stepcounter{utterance}  
    & & \multicolumn{4}{p{0.6\linewidth}}{
        \cellcolor[rgb]{0.9,0.9,0.9}{
            \makecell[{{p{\linewidth}}}]{
                \texttt{\tiny{[P1$\langle$GM]}}
                \texttt{Il{-}messaġġ tiegħek ġie mgħoddi lill{-}plejer l{-}ieħor.} \\
\\ 
\\ 
\texttt{Bħalissa qed tilagħbu r{-}round 12 minn massimu ta' 20.} \\
\texttt{Intom flimkien akkumulajtu 6 minn 12 penali.} \\
\texttt{Il{-}plejer l{-}ieħor bagħat dan il{-}messaġġ:} \\
\\ 
\\ 
\texttt{"Aħna qegħdin ninsa l{-}ħajtar f'(2,1)! Għaxda, l{-}għan li qed nitlobu mhux possibbli. Il{-}kordinati vojta f'Riga 1 huma: (1,1), (3,1), (4,1), (5,1), (6,1). (2,1) u (7,1) huma mblokkati. Għandna 5 oġġetti u 5 kordinati vojta {-} perfett! Nitlobu għan ġdid: W(1,1), I(3,1), T(4,1), C(5,1), H(6,1). Dak huwa l{-}istess li għandi issa! Int wasalt għal W(1,1), T(3,1), C(4,1), H(5,1). Għandek bżonn tċaqlaq I għal (6,1) u T għal (3,1) {-} tista' tagħmel dan? Għidli fejn qegħdin l{-}oġġetti tiegħek ISSA u jekk tista' tagħmel dan."} \\
\\ 
\\ 
\texttt{X'inhu l{-}kmand li jmiss tiegħek?} \\
            }
        }
    }
    & \\ \\

    \theutterance \stepcounter{utterance}  
    & \multicolumn{4}{p{0.6\linewidth}}{
        \cellcolor[rgb]{0.9,0.9,0.9}{
            \makecell[{{p{\linewidth}}}]{
                \texttt{\tiny{[P1$\rangle$GM]}}
                \texttt{GĦID: L{-}oġġetti tiegħi ISSA: W(1,1), I(7,1), T(3,1), C(4,1), H(5,1). T qed jispiċċa f'(3,1) hekk kif rridna. Għandi bżonn: 1) Inħalli (7,1) vojta (mexxi I), 2) Nċaqlaq I għal (6,1). Se nmexxi I minn (7,1) għal (6,1) meta tkun vojta. Issa (6,1) hija vojta. Nagħmel dan issa.} \\
            }
        }
    }
    & & \\ \\

    \theutterance \stepcounter{utterance}  
    & & & \multicolumn{2}{p{0.3\linewidth}}{
        \cellcolor[rgb]{0.9,0.9,0.9}{
            \makecell[{{p{\linewidth}}}]{
                \texttt{\tiny{[GM$|$GM]}}
                \texttt{GĦID: L{-}oġġetti tiegħi ISSA: W(1,1), I(7,1), T(3,1), C(4,1), H(5,1). T qed jispiċċa f'(3,1) hekk kif rridna. Għandi bżonn: 1) Inħalli (7,1) vojta (mexxi I), 2) Nċaqlaq I għal (6,1). Se nmexxi I minn (7,1) għal (6,1) meta tkun vojta. Issa (6,1) hija vojta. Nagħmel dan issa.} \\
            }
        }
    }
    & & \\ \\

    \theutterance \stepcounter{utterance}  
    & & \multicolumn{4}{p{0.6\linewidth}}{
        \cellcolor[rgb]{0.9,0.9,0.9}{
            \makecell[{{p{\linewidth}}}]{
                \texttt{\tiny{[GM$\rangle$P2]}}
                \texttt{Il{-}messaġġ tiegħek ġie mgħoddi lill{-}plejer l{-}ieħor.} \\
\\ 
\\ 
\texttt{Bħalissa qed tilagħbu r{-}round 13 minn massimu ta' 20.} \\
\texttt{Intom flimkien akkumulajtu 6 minn 12 penali.} \\
\texttt{Il{-}plejer l{-}ieħor bagħat dan il{-}messaġġ:} \\
\\ 
\\ 
\texttt{"L{-}oġġetti tiegħi ISSA: W(1,1), I(7,1), T(3,1), C(4,1), H(5,1). T qed jispiċċa f'(3,1) hekk kif rridna. Għandi bżonn: 1) Inħalli (7,1) vojta (mexxi I), 2) Nċaqlaq I għal (6,1). Se nmexxi I minn (7,1) għal (6,1) meta tkun vojta. Issa (6,1) hija vojta. Nagħmel dan issa."} \\
\\ 
\\ 
\texttt{X'inhu l{-}kmand li jmiss tiegħek?} \\
            }
        }
    }
    & \\ \\

    \theutterance \stepcounter{utterance}  
    & & & \multicolumn{4}{p{0.6\linewidth}}{
        \cellcolor[rgb]{0.9,0.9,0.9}{
            \makecell[{{p{\linewidth}}}]{
                \texttt{\tiny{[GM$\langle$P2]}}
                \texttt{GĦID: Perfett! Mexxi 'I' tiegħek għal (6,1). Meta wassal, l{-}oġġetti tiegħek se jkunu: W(1,1), T(3,1), C(4,1), H(5,1), I(6,1). L{-}oġġetti tiegħi huma identici: W(1,1), I(3,1), T(4,1), C(5,1), H(7,1)... الاستنّا، هادي مش متطابقة. نخدمولي: ل{-}هدف ديالي W(1,1), I(3,1), T(4,1), C(5,1), H(6,1). إنت عندك W(1,1), T(3,1), C(4,1), H(5,1), I(6,1). T و C مقلوبين، و H و I مقلوبين. نحتاجو نبدلو توا: إنت تحرك T لـ (4,1) و C لـ (3,1)، و H لـ (6,1) و I لـ (5,1). بس (3,1) و (4,1) مش فاضيين. الحل: إنت تحرك T لـ (6,1) مؤقتاً، تحرك C لـ (3,1)، T لـ (4,1)، H لـ (6,1)، I لـ (5,1). معقد. الأحسن: نتفقوا على هدف جديد يناسب اللوحين. هدف نهائي: W(1,1), I(3,1), T(4,1), C(5,1), H(7,1). إنت عندك H(5,1) و I(6,1). تحرك H لـ (7,1) و I لـ (3,1) {-} بس I عندي (3,1). مبدئياً: خليك إنت تحرك H لـ (7,1) و I لـ (2,1) مش ممكن. خلينا نقولو: الهدف W(1,1), I(3,1), T(4,1), C(5,1), H(6,1). إنت تحرك H من (5,1) لـ (6,1) بس (6,1) فيها I تاعك. تحرك I تاعك لـ (7,1)؟ (7,1) فيها حائط. للا. الحل الوحيد: إنت تحرك H لـ (6,1) و I لـ (7,1) مش ممكن. خلينا نطلبوا الهدف النهائي: W(1,1), I(3,1), T(4,1), C(5,1), H(7,1). إنت تحرك H لـ (7,1) {-} ممكن؟ (7,1) فيها حائط للا. يئس. الهدف الوحيد الممكن: W(1,1), I(3,1), T(4,1), C(5,1), H(6,1). إنت محتاج تحرك H لـ (6,1) و I لـ (5,1). ممكن؟ (5,1) فيها C تاعك، (6,1) فيها I تاعك. تبدلو أماكنهم: إنت تحرك I لـ (7,1) للا. تحرك C لـ (6,1) للا. خليني نفكر... إنت تحرك I لـ (3,1) للا فيها T. توا إنت عندك: W(1,1), T(3,1), C(4,1), H(5,1), I(6,1). جرب: 1. تحرك I من (6,1) لـ (7,1) {-} حائط. 2. تحرك H من (5,1) لـ (6,1) {-} فيها I. 3. تحرك C من (4,1) لـ (5,1) {-} فيها H. 4. تحرك T من (3,1) لـ (4,1) {-} فيها C. 5. تحرك I لـ (3,1) {-} فاضية توا. 6. تحرك T لـ (4,1). 7. تحرك C لـ (5,1). 8. تحرك H لـ (6,1). هكاك نوصلكم للهدف W(1,1), I(3,1), T(4,1), C(5,1), H(6,1). إنت تبدا: تحرك I من (6,1) لـ (3,1) {-} للا مش فاضية، فيها T. تحرك T الاول؟ تحرك T لـ (2,1) {-} حائط. لـ (4,1) {-} فيها C. لـ (3,2)؟ ممكن. خليني نقولك: نتفقوا إن إنت تعمل sequence تاع التحريكات هادي، و نيا نتحرك H تاعيا من (7,1) لـ (6,1) توا (6,1) فاضية عنديا). إنت: 1. T (3,1){-}>(3,2) 2. I (6,1){-}>(3,1) 3. T (3,2){-}>(4,1) 4. C (4,1){-}>(5,1) 5. H (5,1){-}>(6,1). نيا نتحرك H (7,1){-}>(6,1) توا. موافق؟} \\
            }
        }
    }
     \\ \\

    \theutterance \stepcounter{utterance}  
    & & & \multicolumn{2}{p{0.3\linewidth}}{
        \cellcolor[rgb]{0.9,0.9,0.9}{
            \makecell[{{p{\linewidth}}}]{
                \texttt{\tiny{[GM$|$GM]}}
                \texttt{GĦID: Perfett! Mexxi 'I' tiegħek għal (6,1). Meta wassal, l{-}oġġetti tiegħek se jkunu: W(1,1), T(3,1), C(4,1), H(5,1), I(6,1). L{-}oġġetti tiegħi huma identici: W(1,1), I(3,1), T(4,1), C(5,1), H(7,1)... الاستنّا، هادي مش متطابقة. نخدمولي: ل{-}هدف ديالي W(1,1), I(3,1), T(4,1), C(5,1), H(6,1). إنت عندك W(1,1), T(3,1), C(4,1), H(5,1), I(6,1). T و C مقلوبين، و H و I مقلوبين. نحتاجو نبدلو توا: إنت تحرك T لـ (4,1) و C لـ (3,1)، و H لـ (6,1) و I لـ (5,1). بس (3,1) و (4,1) مش فاضيين. الحل: إنت تحرك T لـ (6,1) مؤقتاً، تحرك C لـ (3,1)، T لـ (4,1)، H لـ (6,1)، I لـ (5,1). معقد. الأحسن: نتفقوا على هدف جديد يناسب اللوحين. هدف نهائي: W(1,1), I(3,1), T(4,1), C(5,1), H(7,1). إنت عندك H(5,1) و I(6,1). تحرك H لـ (7,1) و I لـ (3,1) {-} بس I عندي (3,1). مبدئياً: خليك إنت تحرك H لـ (7,1) و I لـ (2,1) مش ممكن. خلينا نقولو: الهدف W(1,1), I(3,1), T(4,1), C(5,1), H(6,1). إنت تحرك H من (5,1) لـ (6,1) بس (6,1) فيها I تاعك. تحرك I تاعك لـ (7,1)؟ (7,1) فيها حائط. للا. الحل الوحيد: إنت تحرك H لـ (6,1) و I لـ (7,1) مش ممكن. خلينا نطلبوا الهدف النهائي: W(1,1), I(3,1), T(4,1), C(5,1), H(7,1). إنت تحرك H لـ (7,1) {-} ممكن؟ (7,1) فيها حائط للا. يئس. الهدف الوحيد الممكن: W(1,1), I(3,1), T(4,1), C(5,1), H(6,1). إنت محتاج تحرك H لـ (6,1) و I لـ (5,1). ممكن؟ (5,1) فيها C تاعك، (6,1) فيها I تاعك. تبدلو أماكنهم: إنت تحرك I لـ (7,1) للا. تحرك C لـ (6,1) للا. خليني نفكر... إنت تحرك I لـ (3,1) للا فيها T. توا إنت عندك: W(1,1), T(3,1), C(4,1), H(5,1), I(6,1). جرب: 1. تحرك I من (6,1) لـ (7,1) {-} حائط. 2. تحرك H من (5,1) لـ (6,1) {-} فيها I. 3. تحرك C من (4,1) لـ (5,1) {-} فيها H. 4. تحرك T من (3,1) لـ (4,1) {-} فيها C. 5. تحرك I لـ (3,1) {-} فاضية توا. 6. تحرك T لـ (4,1). 7. تحرك C لـ (5,1). 8. تحرك H لـ (6,1). هكاك نوصلكم للهدف W(1,1), I(3,1), T(4,1), C(5,1), H(6,1). إنت تبدا: تحرك I من (6,1) لـ (3,1) {-} للا مش فاضية، فيها T. تحرك T الاول؟ تحرك T لـ (2,1) {-} حائط. لـ (4,1) {-} فيها C. لـ (3,2)؟ ممكن. خليني نقولك: نتفقوا إن إنت تعمل sequence تاع التحريكات هادي، و نيا نتحرك H تاعيا من (7,1) لـ (6,1) توا (6,1) فاضية عنديا). إنت: 1. T (3,1){-}>(3,2) 2. I (6,1){-}>(3,1) 3. T (3,2){-}>(4,1) 4. C (4,1){-}>(5,1) 5. H (5,1){-}>(6,1). نيا نتحرك H (7,1){-}>(6,1) توا. موافق؟} \\
            }
        }
    }
    & & \\ \\

    \theutterance \stepcounter{utterance}  
    & & \multicolumn{4}{p{0.6\linewidth}}{
        \cellcolor[rgb]{0.9,0.9,0.9}{
            \makecell[{{p{\linewidth}}}]{
                \texttt{\tiny{[P1$\langle$GM]}}
                \texttt{Il{-}messaġġ tiegħek ġie mgħoddi lill{-}plejer l{-}ieħor.} \\
\\ 
\\ 
\texttt{Bħalissa qed tilagħbu r{-}round 13 minn massimu ta' 20.} \\
\texttt{Intom flimkien akkumulajtu 6 minn 12 penali.} \\
\texttt{Il{-}plejer l{-}ieħor bagħat dan il{-}messaġġ:} \\
\\ 
\\ 
\texttt{"Perfett! Mexxi 'I' tiegħek għal (6,1). Meta wassal, l{-}oġġetti tiegħek se jkunu: W(1,1), T(3,1), C(4,1), H(5,1), I(6,1). L{-}oġġetti tiegħi huma identici: W(1,1), I(3,1), T(4,1), C(5,1), H(7,1)... الاستنّا، هادي مش متطابقة. نخدمولي: ل{-}هدف ديالي W(1,1), I(3,1), T(4,1), C(5,1), H(6,1). إنت عندك W(1,1), T(3,1), C(4,1), H(5,1), I(6,1). T و C مقلوبين، و H و I مقلوبين. نحتاجو نبدلو توا: إنت تحرك T لـ (4,1) و C لـ (3,1)، و H لـ (6,1) و I لـ (5,1). بس (3,1) و (4,1) مش فاضيين. الحل: إنت تحرك T لـ (6,1) مؤقتاً، تحرك C لـ (3,1)، T لـ (4,1)، H لـ (6,1)، I لـ (5,1). معقد. الأحسن: نتفقوا على هدف جديد يناسب اللوحين. هدف نهائي: W(1,1), I(3,1), T(4,1), C(5,1), H(7,1). إنت عندك H(5,1) و I(6,1). تحرك H لـ (7,1) و I لـ (3,1) {-} بس I عندي (3,1). مبدئياً: خليك إنت تحرك H لـ (7,1) و I لـ (2,1) مش ممكن. خلينا نقولو: الهدف W(1,1), I(3,1), T(4,1), C(5,1), H(6,1). إنت تحرك H من (5,1) لـ (6,1) بس (6,1) فيها I تاعك. تحرك I تاعك لـ (7,1)؟ (7,1) فيها حائط. للا. الحل الوحيد: إنت تحرك H لـ (6,1) و I لـ (7,1) مش ممكن. خلينا نطلبوا الهدف النهائي: W(1,1), I(3,1), T(4,1), C(5,1), H(7,1). إنت تحرك H لـ (7,1) {-} ممكن؟ (7,1) فيها حائط للا. يئس. الهدف الوحيد الممكن: W(1,1), I(3,1), T(4,1), C(5,1), H(6,1). إنت محتاج تحرك H لـ (6,1) و I لـ (5,1). ممكن؟ (5,1) فيها C تاعك، (6,1) فيها I تاعك. تبدلو أماكنهم: إنت تحرك I لـ (7,1) للا. تحرك C لـ (6,1) للا. خليني نفكر... إنت تحرك I لـ (3,1) للا فيها T. توا إنت عندك: W(1,1), T(3,1), C(4,1), H(5,1), I(6,1). جرب: 1. تحرك I من (6,1) لـ (7,1) {-} حائط. 2. تحرك H من (5,1) لـ (6,1) {-} فيها I. 3. تحرك C من (4,1) لـ (5,1) {-} فيها H. 4. تحرك T من (3,1) لـ (4,1) {-} فيها C. 5. تحرك I لـ (3,1) {-} فاضية توا. 6. تحرك T لـ (4,1). 7. تحرك C لـ (5,1). 8. تحرك H لـ (6,1). هكاك نوصلكم للهدف W(1,1), I(3,1), T(4,1), C(5,1), H(6,1). إنت تبدا: تحرك I من (6,1) لـ (3,1) {-} للا مش فاضية، فيها T. تحرك T الاول؟ تحرك T لـ (2,1) {-} حائط. لـ (4,1) {-} فيها C. لـ (3,2)؟ ممكن. خليني نقولك: نتفقوا إن إنت تعمل sequence تاع التحريكات هادي، و نيا نتحرك H تاعيا من (7,1) لـ (6,1) توا (6,1) فاضية عنديا). إنت: 1. T (3,1){-}>(3,2) 2. I (6,1){-}>(3,1) 3. T (3,2){-}>(4,1) 4. C (4,1){-}>(5,1) 5. H (5,1){-}>(6,1). نيا نتحرك H (7,1){-}>(6,1) توا. موافق؟"} \\
\\ 
\\ 
\texttt{X'inhu l{-}kmand li jmiss tiegħek?} \\
            }
        }
    }
    & \\ \\

    \theutterance \stepcounter{utterance}  
    & \multicolumn{4}{p{0.6\linewidth}}{
        \cellcolor[rgb]{0.9,0.9,0.9}{
            \makecell[{{p{\linewidth}}}]{
                \texttt{\tiny{[P1$\rangle$GM]}}
                \texttt{GĦID: Moqbel nipproċedura tiegħek. Inbidu: T minn (3,1) għal (3,2).} \\
            }
        }
    }
    & & \\ \\

    \theutterance \stepcounter{utterance}  
    & & & \multicolumn{2}{p{0.3\linewidth}}{
        \cellcolor[rgb]{0.9,0.9,0.9}{
            \makecell[{{p{\linewidth}}}]{
                \texttt{\tiny{[GM$|$GM]}}
                \texttt{GĦID: Moqbel nipproċedura tiegħek. Inbidu: T minn (3,1) għal (3,2).} \\
            }
        }
    }
    & & \\ \\

    \theutterance \stepcounter{utterance}  
    & & \multicolumn{4}{p{0.6\linewidth}}{
        \cellcolor[rgb]{0.9,0.9,0.9}{
            \makecell[{{p{\linewidth}}}]{
                \texttt{\tiny{[GM$\rangle$P2]}}
                \texttt{Il{-}messaġġ tiegħek ġie mgħoddi lill{-}plejer l{-}ieħor.} \\
\\ 
\\ 
\texttt{Bħalissa qed tilagħbu r{-}round 14 minn massimu ta' 20.} \\
\texttt{Intom flimkien akkumulajtu 6 minn 12 penali.} \\
\texttt{Il{-}plejer l{-}ieħor bagħat dan il{-}messaġġ:} \\
\\ 
\\ 
\texttt{"Moqbel nipproċedura tiegħek. Inbidu: T minn (3,1) għal (3,2)."} \\
\\ 
\\ 
\texttt{X'inhu l{-}kmand li jmiss tiegħek?} \\
            }
        }
    }
    & \\ \\

    \theutterance \stepcounter{utterance}  
    & & & \multicolumn{4}{p{0.6\linewidth}}{
        \cellcolor[rgb]{0.9,0.9,0.9}{
            \makecell[{{p{\linewidth}}}]{
                \texttt{\tiny{[GM$\langle$P2]}}
                \texttt{ĊAQLEQ: H, (6, 1)} \\
            }
        }
    }
     \\ \\

    \theutterance \stepcounter{utterance}  
    & & & \multicolumn{2}{p{0.3\linewidth}}{
        \cellcolor[rgb]{0.9,0.9,0.9}{
            \makecell[{{p{\linewidth}}}]{
                \texttt{\tiny{[GM$|$GM]}}
                \texttt{ĊAQLEQ: H, (6, 1)} \\
            }
        }
    }
    & & \\ \\

    \theutterance \stepcounter{utterance}  
    & & & \multicolumn{2}{p{0.3\linewidth}}{
        \cellcolor[rgb]{0.9,0.9,0.9}{
            \makecell[{{p{\linewidth}}}]{
                \texttt{\tiny{[GM$|$GM]}}
                \texttt{Penali: (6,1) mhijiex vojta, iżda fiha '│'.} \\
\\ 
\texttt{Intom flimkien akkumulajtu 7 minn 12 penali.} \\
\texttt{Jekk jogħġbok erġa' pprova!} \\
            }
        }
    }
    & & \\ \\

    \theutterance \stepcounter{utterance}  
    & & \multicolumn{4}{p{0.6\linewidth}}{
        \cellcolor[rgb]{0.9,0.9,0.9}{
            \makecell[{{p{\linewidth}}}]{
                \texttt{\tiny{[GM$\rangle$P2]}}
                \texttt{Penali: (6,1) mhijiex vojta, iżda fiha '│'.} \\
\\ 
\texttt{Intom flimkien akkumulajtu 7 minn 12 penali.} \\
\texttt{Jekk jogħġbok erġa' pprova!} \\
            }
        }
    }
    & \\ \\

    \theutterance \stepcounter{utterance}  
    & & & \multicolumn{4}{p{0.6\linewidth}}{
        \cellcolor[rgb]{0.9,0.9,0.9}{
            \makecell[{{p{\linewidth}}}]{
                \texttt{\tiny{[GM$\langle$P2]}}
                \texttt{GĦID: L{-}ħajta vertikali f'(6,1) qed tbloċċja l{-}pass ta' H. Għandna bżonn nċaqlaq 'I' minn (3,1) għal (2,1) għal ma niftaħx il{-}pass għal H. Jekk inti testenniex 'I', nistgħu nkompletu.} \\
            }
        }
    }
     \\ \\

    \theutterance \stepcounter{utterance}  
    & & & \multicolumn{2}{p{0.3\linewidth}}{
        \cellcolor[rgb]{0.9,0.9,0.9}{
            \makecell[{{p{\linewidth}}}]{
                \texttt{\tiny{[GM$|$GM]}}
                \texttt{GĦID: L{-}ħajta vertikali f'(6,1) qed tbloċċja l{-}pass ta' H. Għandna bżonn nċaqlaq 'I' minn (3,1) għal (2,1) għal ma niftaħx il{-}pass għal H. Jekk inti testenniex 'I', nistgħu nkompletu.} \\
            }
        }
    }
    & & \\ \\

    \theutterance \stepcounter{utterance}  
    & & \multicolumn{4}{p{0.6\linewidth}}{
        \cellcolor[rgb]{0.9,0.9,0.9}{
            \makecell[{{p{\linewidth}}}]{
                \texttt{\tiny{[P1$\langle$GM]}}
                \texttt{Il{-}messaġġ tiegħek ġie mgħoddi lill{-}plejer l{-}ieħor.} \\
\\ 
\\ 
\texttt{Bħalissa qed tilagħbu r{-}round 14 minn massimu ta' 20.} \\
\texttt{Intom flimkien akkumulajtu 7 minn 12 penali.} \\
\texttt{Il{-}plejer l{-}ieħor bagħat dan il{-}messaġġ:} \\
\\ 
\\ 
\texttt{"L{-}ħajta vertikali f'(6,1) qed tbloċċja l{-}pass ta' H. Għandna bżonn nċaqlaq 'I' minn (3,1) għal (2,1) għal ma niftaħx il{-}pass għal H. Jekk inti testenniex 'I', nistgħu nkompletu."} \\
\\ 
\\ 
\texttt{X'inhu l{-}kmand li jmiss tiegħek?} \\
            }
        }
    }
    & \\ \\

    \theutterance \stepcounter{utterance}  
    & \multicolumn{4}{p{0.6\linewidth}}{
        \cellcolor[rgb]{0.9,0.9,0.9}{
            \makecell[{{p{\linewidth}}}]{
                \texttt{\tiny{[P1$\rangle$GM]}}
                \texttt{ĊAQLEQ: T, (3, 2)} \\
            }
        }
    }
    & & \\ \\

    \theutterance \stepcounter{utterance}  
    & & & \multicolumn{2}{p{0.3\linewidth}}{
        \cellcolor[rgb]{0.9,0.9,0.9}{
            \makecell[{{p{\linewidth}}}]{
                \texttt{\tiny{[GM$|$GM]}}
                \texttt{ĊAQLEQ: T, (3, 2)} \\
            }
        }
    }
    & & \\ \\

    \theutterance \stepcounter{utterance}  
    & & & \multicolumn{2}{p{0.3\linewidth}}{
        \cellcolor[rgb]{0.9,0.9,0.9}{
            \makecell[{{p{\linewidth}}}]{
                \texttt{\tiny{[GM$|$GM]}}
                \texttt{Penali: (3,2) mhijiex vojta, iżda fiha '─'.} \\
\\ 
\texttt{Intom flimkien akkumulajtu 8 minn 12 penali.} \\
\texttt{Jekk jogħġbok erġa' pprova!} \\
            }
        }
    }
    & & \\ \\

    \theutterance \stepcounter{utterance}  
    & & \multicolumn{4}{p{0.6\linewidth}}{
        \cellcolor[rgb]{0.9,0.9,0.9}{
            \makecell[{{p{\linewidth}}}]{
                \texttt{\tiny{[P1$\langle$GM]}}
                \texttt{Penali: (3,2) mhijiex vojta, iżda fiha '─'.} \\
\\ 
\texttt{Intom flimkien akkumulajtu 8 minn 12 penali.} \\
\texttt{Jekk jogħġbok erġa' pprova!} \\
            }
        }
    }
    & \\ \\

    \theutterance \stepcounter{utterance}  
    & \multicolumn{4}{p{0.6\linewidth}}{
        \cellcolor[rgb]{0.9,0.9,0.9}{
            \makecell[{{p{\linewidth}}}]{
                \texttt{\tiny{[P1$\rangle$GM]}}
                \texttt{GĦID: Penali għax (3,2) mhix vojta (fih '─'). Fejn huma l{-}ċelloli vojta madwar (3,1)? Qed nara li (3,1) mħux mħallt minn ħijtar. X'pożizzjonijiet vojta għandek int?} \\
            }
        }
    }
    & & \\ \\

    \theutterance \stepcounter{utterance}  
    & & & \multicolumn{2}{p{0.3\linewidth}}{
        \cellcolor[rgb]{0.9,0.9,0.9}{
            \makecell[{{p{\linewidth}}}]{
                \texttt{\tiny{[GM$|$GM]}}
                \texttt{GĦID: Penali għax (3,2) mhix vojta (fih '─'). Fejn huma l{-}ċelloli vojta madwar (3,1)? Qed nara li (3,1) mħux mħallt minn ħijtar. X'pożizzjonijiet vojta għandek int?} \\
            }
        }
    }
    & & \\ \\

    \theutterance \stepcounter{utterance}  
    & & \multicolumn{4}{p{0.6\linewidth}}{
        \cellcolor[rgb]{0.9,0.9,0.9}{
            \makecell[{{p{\linewidth}}}]{
                \texttt{\tiny{[GM$\rangle$P2]}}
                \texttt{Il{-}messaġġ tiegħek ġie mgħoddi lill{-}plejer l{-}ieħor.} \\
\\ 
\\ 
\texttt{Bħalissa qed tilagħbu r{-}round 15 minn massimu ta' 20.} \\
\texttt{Intom flimkien akkumulajtu 8 minn 12 penali.} \\
\texttt{Il{-}plejer l{-}ieħor bagħat dan il{-}messaġġ:} \\
\\ 
\\ 
\texttt{"Penali għax (3,2) mhix vojta (fih '─'). Fejn huma l{-}ċelloli vojta madwar (3,1)? Qed nara li (3,1) mħux mħallt minn ħijtar. X'pożizzjonijiet vojta għandek int?"} \\
\\ 
\\ 
\texttt{X'inhu l{-}kmand li jmiss tiegħek?} \\
            }
        }
    }
    & \\ \\

    \theutterance \stepcounter{utterance}  
    & & & \multicolumn{4}{p{0.6\linewidth}}{
        \cellcolor[rgb]{0.9,0.9,0.9}{
            \makecell[{{p{\linewidth}}}]{
                \texttt{\tiny{[GM$\langle$P2]}}
                \texttt{GĦID: Ċelloli vojta madwar (3,1) f'grilja tiegħi: (3,2) mhix vojta (fih '─'), (2,1) mhix vojta (fih '│'), (4,1) vojta, (3,0) barra mil{-}grilja. Għalhekk, T f'(3,1) mhux movibbli f'daqqika. } \\
\\ 
\texttt{L{-}oġġetti tiegħi li huma movibbli: I f'(3,1) {-}> tista' tmur (4,1); C f'(5,1) {-}> tista' tmur (4,1) jew (6,1); H f'(7,1) {-}> mhux movibbli (ħajt f'(6,1)). } \\
\\ 
\texttt{L{-}oġġetti tiegħek: T f'(3,1) mhux movibbli; I f'(7,1) {-}> tista' tmur (6,1) jew (7,2); C f'(4,1) {-}> tista' tmur (3,1) mhux vojta, (5,1) mhux vojta, (4,2); H f'(5,1) {-}> tista' tmur (4,1) mhux vojta, (6,1), (5,2).} \\
\\ 
\texttt{Nista' nċaqlaq 'I' tiegħi minn (3,1) għal (4,1) meta tkun vojta. Tista' inti tċaqlaq 'H' tiegħek minn (5,1) għal (6,1) issa? (6,1) hija vojta f'grilja tiegħek?} \\
            }
        }
    }
     \\ \\

    \theutterance \stepcounter{utterance}  
    & & & \multicolumn{2}{p{0.3\linewidth}}{
        \cellcolor[rgb]{0.9,0.9,0.9}{
            \makecell[{{p{\linewidth}}}]{
                \texttt{\tiny{[GM$|$GM]}}
                \texttt{GĦID: Ċelloli vojta madwar (3,1) f'grilja tiegħi: (3,2) mhix vojta (fih '─'), (2,1) mhix vojta (fih '│'), (4,1) vojta, (3,0) barra mil{-}grilja. Għalhekk, T f'(3,1) mhux movibbli f'daqqika. } \\
\\ 
\texttt{L{-}oġġetti tiegħi li huma movibbli: I f'(3,1) {-}> tista' tmur (4,1); C f'(5,1) {-}> tista' tmur (4,1) jew (6,1); H f'(7,1) {-}> mhux movibbli (ħajt f'(6,1)). } \\
\\ 
\texttt{L{-}oġġetti tiegħek: T f'(3,1) mhux movibbli; I f'(7,1) {-}> tista' tmur (6,1) jew (7,2); C f'(4,1) {-}> tista' tmur (3,1) mhux vojta, (5,1) mhux vojta, (4,2); H f'(5,1) {-}> tista' tmur (4,1) mhux vojta, (6,1), (5,2).} \\
\\ 
\texttt{Nista' nċaqlaq 'I' tiegħi minn (3,1) għal (4,1) meta tkun vojta. Tista' inti tċaqlaq 'H' tiegħek minn (5,1) għal (6,1) issa? (6,1) hija vojta f'grilja tiegħek?} \\
            }
        }
    }
    & & \\ \\

    \theutterance \stepcounter{utterance}  
    & & & \multicolumn{2}{p{0.3\linewidth}}{
        \cellcolor[rgb]{0.9,0.9,0.9}{
            \makecell[{{p{\linewidth}}}]{
                \texttt{\tiny{[GM$|$GM]}}
                \texttt{Invalid format: tail is not empty: '} \\
\\ 
\texttt{L{-}oġġetti tiegħi li huma movibbli: I f'(3,1) {-}> tista' tmur (4,1); C f'(5,1) {-}> tista' tmur (4,1) jew (6,1); H f'(7,1) {-}> mhux movibbli (ħajt f'(6,1)). } \\
\\ 
\texttt{L{-}oġġetti tiegħek: T f'(3,1) mhux movibbli; I f'(7,1) {-}> tista' tmur (6,1) jew (7,2); C f'(4,1) {-}> tista' tmur (3,1) mhux vojta, (5,1) mhux vojta, (4,2); H f'(5,1) {-}> tista' tmur (4,1) mhux vojta, (6,1), (5,2).} \\
\\ 
\texttt{Nista' nċaqlaq 'I' tiegħi minn (3,1) għal (4,1) meta tkun vojta. Tista' inti tċaqlaq 'H' tiegħek minn (5,1) għal (6,1) issa? (6,1) hija vojta f'grilja tiegħek?} \\
            }
        }
    }
    & & \\ \\

    \theutterance \stepcounter{utterance}  
    & & \multicolumn{4}{p{0.6\linewidth}}{
        \cellcolor[rgb]{0.9,0.9,0.9}{
            \makecell[{{p{\linewidth}}}]{
                \texttt{\tiny{[GM$\rangle$P2]}}
                \texttt{Penali: Il{-}messaġġ tiegħek m'għandux ikun fih xejn wara l{-}kmand!} \\
\\ 
\texttt{Kun żgur li r{-}risposta tiegħek fiha biss jew `GĦID: <MESSAĠĠ>` jew `ĊAQLEQ: <OĠĠETT>, (<X>, <Y>)`, u xejn ieħor!} \\
\\ 
\texttt{Intom flimkien akkumulajtu 9 minn 12 penali.} \\
\texttt{Jekk jogħġbok erġa' pprova!} \\
            }
        }
    }
    & \\ \\

    \theutterance \stepcounter{utterance}  
    & & & \multicolumn{4}{p{0.6\linewidth}}{
        \cellcolor[rgb]{0.9,0.9,0.9}{
            \makecell[{{p{\linewidth}}}]{
                \texttt{\tiny{[GM$\langle$P2]}}
                \texttt{GĦID: Ċelloli vojta madwar (3,1) f'grilja tiegħi: (3,2) mhix vojta (fih '─'), (2,1) mhix vojta (fih '│'), (4,1) vojta, (3,0) barra mil{-}grilja. Għalhekk, T f'(3,1) mhux movibbli f'daqqika. L{-}oġġetti tiegħi li huma movibbli: I f'(3,1) {-}> tista' tmur (4,1); C f'(5,1) {-}> tista' tmur (4,1) jew (6,1); H f'(7,1) {-}> mhux movibbli (ħajt f'(6,1)). L{-}oġġetti tiegħek: T f'(3,1) mhux movibbli; I f'(7,1) {-}> tista' tmur (6,1) jew (7,2); C f'(4,1) {-}> tista' tmur (3,1) mhux vojta, (5,1) mhux vojta, (4,2); H f'(5,1) {-}> tista' tmur (4,1) mhux vojta, (6,1), (5,2). Nista' nċaqlaq 'I' tiegħi minn (3,1) għal (4,1) meta tkun vojta. Tista' inti tċaqlaq 'H' tiegħek minn (5,1) għal (6,1) issa? (6,1) hija vojta f'grilja tiegħek?} \\
            }
        }
    }
     \\ \\

    \theutterance \stepcounter{utterance}  
    & & & \multicolumn{2}{p{0.3\linewidth}}{
        \cellcolor[rgb]{0.9,0.9,0.9}{
            \makecell[{{p{\linewidth}}}]{
                \texttt{\tiny{[GM$|$GM]}}
                \texttt{GĦID: Ċelloli vojta madwar (3,1) f'grilja tiegħi: (3,2) mhix vojta (fih '─'), (2,1) mhix vojta (fih '│'), (4,1) vojta, (3,0) barra mil{-}grilja. Għalhekk, T f'(3,1) mhux movibbli f'daqqika. L{-}oġġetti tiegħi li huma movibbli: I f'(3,1) {-}> tista' tmur (4,1); C f'(5,1) {-}> tista' tmur (4,1) jew (6,1); H f'(7,1) {-}> mhux movibbli (ħajt f'(6,1)). L{-}oġġetti tiegħek: T f'(3,1) mhux movibbli; I f'(7,1) {-}> tista' tmur (6,1) jew (7,2); C f'(4,1) {-}> tista' tmur (3,1) mhux vojta, (5,1) mhux vojta, (4,2); H f'(5,1) {-}> tista' tmur (4,1) mhux vojta, (6,1), (5,2). Nista' nċaqlaq 'I' tiegħi minn (3,1) għal (4,1) meta tkun vojta. Tista' inti tċaqlaq 'H' tiegħek minn (5,1) għal (6,1) issa? (6,1) hija vojta f'grilja tiegħek?} \\
            }
        }
    }
    & & \\ \\

    \theutterance \stepcounter{utterance}  
    & & \multicolumn{4}{p{0.6\linewidth}}{
        \cellcolor[rgb]{0.9,0.9,0.9}{
            \makecell[{{p{\linewidth}}}]{
                \texttt{\tiny{[P1$\langle$GM]}}
                \texttt{Il{-}messaġġ tiegħek ġie mgħoddi lill{-}plejer l{-}ieħor.} \\
\\ 
\\ 
\texttt{Bħalissa qed tilagħbu r{-}round 15 minn massimu ta' 20.} \\
\texttt{Intom flimkien akkumulajtu 9 minn 12 penali.} \\
\texttt{Il{-}plejer l{-}ieħor bagħat dan il{-}messaġġ:} \\
\\ 
\\ 
\texttt{"Ċelloli vojta madwar (3,1) f'grilja tiegħi: (3,2) mhix vojta (fih '─'), (2,1) mhix vojta (fih '│'), (4,1) vojta, (3,0) barra mil{-}grilja. Għalhekk, T f'(3,1) mhux movibbli f'daqqika. L{-}oġġetti tiegħi li huma movibbli: I f'(3,1) {-}> tista' tmur (4,1); C f'(5,1) {-}> tista' tmur (4,1) jew (6,1); H f'(7,1) {-}> mhux movibbli (ħajt f'(6,1)). L{-}oġġetti tiegħek: T f'(3,1) mhux movibbli; I f'(7,1) {-}> tista' tmur (6,1) jew (7,2); C f'(4,1) {-}> tista' tmur (3,1) mhux vojta, (5,1) mhux vojta, (4,2); H f'(5,1) {-}> tista' tmur (4,1) mhux vojta, (6,1), (5,2). Nista' nċaqlaq 'I' tiegħi minn (3,1) għal (4,1) meta tkun vojta. Tista' inti tċaqlaq 'H' tiegħek minn (5,1) għal (6,1) issa? (6,1) hija vojta f'grilja tiegħek?"} \\
\\ 
\\ 
\texttt{X'inhu l{-}kmand li jmiss tiegħek?} \\
            }
        }
    }
    & \\ \\

    \theutterance \stepcounter{utterance}  
    & \multicolumn{4}{p{0.6\linewidth}}{
        \cellcolor[rgb]{0.9,0.9,0.9}{
            \makecell[{{p{\linewidth}}}]{
                \texttt{\tiny{[P1$\rangle$GM]}}
                \texttt{GĦID: Iva, (6,1) hija vojta f'grilja tiegħi. Se nċaqlaq 'H' tiegħi minn (5,1) għal (6,1) issa. Dan jiftaħ (5,1) għall{-}passi ġejjin.} \\
            }
        }
    }
    & & \\ \\

    \theutterance \stepcounter{utterance}  
    & & & \multicolumn{2}{p{0.3\linewidth}}{
        \cellcolor[rgb]{0.9,0.9,0.9}{
            \makecell[{{p{\linewidth}}}]{
                \texttt{\tiny{[GM$|$GM]}}
                \texttt{GĦID: Iva, (6,1) hija vojta f'grilja tiegħi. Se nċaqlaq 'H' tiegħi minn (5,1) għal (6,1) issa. Dan jiftaħ (5,1) għall{-}passi ġejjin.} \\
            }
        }
    }
    & & \\ \\

    \theutterance \stepcounter{utterance}  
    & & \multicolumn{4}{p{0.6\linewidth}}{
        \cellcolor[rgb]{0.9,0.9,0.9}{
            \makecell[{{p{\linewidth}}}]{
                \texttt{\tiny{[GM$\rangle$P2]}}
                \texttt{Il{-}messaġġ tiegħek ġie mgħoddi lill{-}plejer l{-}ieħor.} \\
\\ 
\\ 
\texttt{Bħalissa qed tilagħbu r{-}round 16 minn massimu ta' 20.} \\
\texttt{Intom flimkien akkumulajtu 9 minn 12 penali.} \\
\texttt{Il{-}plejer l{-}ieħor bagħat dan il{-}messaġġ:} \\
\\ 
\\ 
\texttt{"Iva, (6,1) hija vojta f'grilja tiegħi. Se nċaqlaq 'H' tiegħi minn (5,1) għal (6,1) issa. Dan jiftaħ (5,1) għall{-}passi ġejjin."} \\
\\ 
\\ 
\texttt{X'inhu l{-}kmand li jmiss tiegħek?} \\
            }
        }
    }
    & \\ \\

    \theutterance \stepcounter{utterance}  
    & & & \multicolumn{4}{p{0.6\linewidth}}{
        \cellcolor[rgb]{0.9,0.9,0.9}{
            \makecell[{{p{\linewidth}}}]{
                \texttt{\tiny{[GM$\langle$P2]}}
                \texttt{GĦID: Perfett. Meta tmexxi 'H' tiegħek għal (6,1), (5,1) tikber vojta f'grilja tiegħek. Wara dak, nista' nċaqlaq 'C' tiegħi minn (5,1) għal (4,1) jew (6,1) f'grilja tiegħi. Int tista' wkoll tċaqlaq 'I' tiegħek minn (7,1) għal (5,1) meta tkun vojta. Nistenna li tagħmel il{-}moviment tiegħek u għidli meta tmexxi 'H' tiegħek.} \\
            }
        }
    }
     \\ \\

    \theutterance \stepcounter{utterance}  
    & & & \multicolumn{2}{p{0.3\linewidth}}{
        \cellcolor[rgb]{0.9,0.9,0.9}{
            \makecell[{{p{\linewidth}}}]{
                \texttt{\tiny{[GM$|$GM]}}
                \texttt{GĦID: Perfett. Meta tmexxi 'H' tiegħek għal (6,1), (5,1) tikber vojta f'grilja tiegħek. Wara dak, nista' nċaqlaq 'C' tiegħi minn (5,1) għal (4,1) jew (6,1) f'grilja tiegħi. Int tista' wkoll tċaqlaq 'I' tiegħek minn (7,1) għal (5,1) meta tkun vojta. Nistenna li tagħmel il{-}moviment tiegħek u għidli meta tmexxi 'H' tiegħek.} \\
            }
        }
    }
    & & \\ \\

    \theutterance \stepcounter{utterance}  
    & & \multicolumn{4}{p{0.6\linewidth}}{
        \cellcolor[rgb]{0.9,0.9,0.9}{
            \makecell[{{p{\linewidth}}}]{
                \texttt{\tiny{[P1$\langle$GM]}}
                \texttt{Il{-}messaġġ tiegħek ġie mgħoddi lill{-}plejer l{-}ieħor.} \\
\\ 
\\ 
\texttt{Bħalissa qed tilagħbu r{-}round 16 minn massimu ta' 20.} \\
\texttt{Intom flimkien akkumulajtu 9 minn 12 penali.} \\
\texttt{Il{-}plejer l{-}ieħor bagħat dan il{-}messaġġ:} \\
\\ 
\\ 
\texttt{"Perfett. Meta tmexxi 'H' tiegħek għal (6,1), (5,1) tikber vojta f'grilja tiegħek. Wara dak, nista' nċaqlaq 'C' tiegħi minn (5,1) għal (4,1) jew (6,1) f'grilja tiegħi. Int tista' wkoll tċaqlaq 'I' tiegħek minn (7,1) għal (5,1) meta tkun vojta. Nistenna li tagħmel il{-}moviment tiegħek u għidli meta tmexxi 'H' tiegħek."} \\
\\ 
\\ 
\texttt{X'inhu l{-}kmand li jmiss tiegħek?} \\
            }
        }
    }
    & \\ \\

    \theutterance \stepcounter{utterance}  
    & \multicolumn{4}{p{0.6\linewidth}}{
        \cellcolor[rgb]{0.9,0.9,0.9}{
            \makecell[{{p{\linewidth}}}]{
                \texttt{\tiny{[P1$\rangle$GM]}}
                \texttt{ĊAQLEQ: H, (6, 1)} \\
            }
        }
    }
    & & \\ \\

    \theutterance \stepcounter{utterance}  
    & & & \multicolumn{2}{p{0.3\linewidth}}{
        \cellcolor[rgb]{0.9,0.9,0.9}{
            \makecell[{{p{\linewidth}}}]{
                \texttt{\tiny{[GM$|$GM]}}
                \texttt{ĊAQLEQ: H, (6, 1)} \\
            }
        }
    }
    & & \\ \\

    \theutterance \stepcounter{utterance}  
    & & & \multicolumn{2}{p{0.3\linewidth}}{
        \cellcolor[rgb]{0.9,0.9,0.9}{
            \makecell[{{p{\linewidth}}}]{
                \texttt{\tiny{[GM$|$GM]}}
                \texttt{Penali: (6,1) mhijiex vojta, iżda fiha '│'.} \\
\\ 
\texttt{Intom flimkien akkumulajtu 10 minn 12 penali.} \\
\texttt{Jekk jogħġbok erġa' pprova!} \\
            }
        }
    }
    & & \\ \\

    \theutterance \stepcounter{utterance}  
    & & \multicolumn{4}{p{0.6\linewidth}}{
        \cellcolor[rgb]{0.9,0.9,0.9}{
            \makecell[{{p{\linewidth}}}]{
                \texttt{\tiny{[P1$\langle$GM]}}
                \texttt{Penali: (6,1) mhijiex vojta, iżda fiha '│'.} \\
\\ 
\texttt{Intom flimkien akkumulajtu 10 minn 12 penali.} \\
\texttt{Jekk jogħġbok erġa' pprova!} \\
            }
        }
    }
    & \\ \\

    \theutterance \stepcounter{utterance}  
    & \multicolumn{4}{p{0.6\linewidth}}{
        \cellcolor[rgb]{0.9,0.9,0.9}{
            \makecell[{{p{\linewidth}}}]{
                \texttt{\tiny{[P1$\rangle$GM]}}
                \texttt{GĦID: (6,1) mhix vojta f'grilja tiegħi (fih '│'). Għalhekk ma nistax nċaqlaq 'H' hemm. Il{-}ħajtar qegħdin jibblokkjaw l{-}moviment. X'pożizzjonijiet vojta għandek int f'Riga 1 u madwarha? Għandna bżonn nkunu realistici dwar fejn nistgħu nmexxu l{-}oġġetti.} \\
            }
        }
    }
    & & \\ \\

    \theutterance \stepcounter{utterance}  
    & & & \multicolumn{2}{p{0.3\linewidth}}{
        \cellcolor[rgb]{0.9,0.9,0.9}{
            \makecell[{{p{\linewidth}}}]{
                \texttt{\tiny{[GM$|$GM]}}
                \texttt{GĦID: (6,1) mhix vojta f'grilja tiegħi (fih '│'). Għalhekk ma nistax nċaqlaq 'H' hemm. Il{-}ħajtar qegħdin jibblokkjaw l{-}moviment. X'pożizzjonijiet vojta għandek int f'Riga 1 u madwarha? Għandna bżonn nkunu realistici dwar fejn nistgħu nmexxu l{-}oġġetti.} \\
            }
        }
    }
    & & \\ \\

    \theutterance \stepcounter{utterance}  
    & & \multicolumn{4}{p{0.6\linewidth}}{
        \cellcolor[rgb]{0.9,0.9,0.9}{
            \makecell[{{p{\linewidth}}}]{
                \texttt{\tiny{[GM$\rangle$P2]}}
                \texttt{Il{-}messaġġ tiegħek ġie mgħoddi lill{-}plejer l{-}ieħor.} \\
\\ 
\\ 
\texttt{Bħalissa qed tilagħbu r{-}round 17 minn massimu ta' 20.} \\
\texttt{Intom flimkien akkumulajtu 10 minn 12 penali.} \\
\texttt{Il{-}plejer l{-}ieħor bagħat dan il{-}messaġġ:} \\
\\ 
\\ 
\texttt{"(6,1) mhix vojta f'grilja tiegħi (fih '│'). Għalhekk ma nistax nċaqlaq 'H' hemm. Il{-}ħajtar qegħdin jibblokkjaw l{-}moviment. X'pożizzjonijiet vojta għandek int f'Riga 1 u madwarha? Għandna bżonn nkunu realistici dwar fejn nistgħu nmexxu l{-}oġġetti."} \\
\\ 
\\ 
\texttt{X'inhu l{-}kmand li jmiss tiegħek?} \\
            }
        }
    }
    & \\ \\

    \theutterance \stepcounter{utterance}  
    & & & \multicolumn{4}{p{0.6\linewidth}}{
        \cellcolor[rgb]{0.9,0.9,0.9}{
            \makecell[{{p{\linewidth}}}]{
                \texttt{\tiny{[GM$\langle$P2]}}
                \texttt{GĦID: F'Riga 1 tiegħi, l{-}ċelloli vojta huma: (1,1) W, (3,1) I, (4,1) T, (5,1) C, (7,1) H. L{-}ċelloli mblokkati mill{-}ħajtar huma: (2,1) '│', (6,1) '│'. Madwar l{-}oġġetti tiegħi: I f'(3,1) tista' tmur (4,1); T f'(4,1) mhux movibbli; C f'(5,1) tista' tmur (4,1) jew (7,1) mhux vojta; H f'(7,1) mhux movibbli. } \\
\\ 
\texttt{Għalhekk, il{-}movimenti possibbli f'Riga 1 huma limitati. Nista' nċaqlaq 'I' minn (3,1) għal (4,1) meta tkun vojta. } \\
\\ 
\texttt{L{-}ebda oġġett ma jista' jidħol f'(2,1) jew (6,1) għax huma ħajtar permanenti. } \\
\\ 
\texttt{Għandna bżonn niftakru għan realistiku li jikseb pożizzjonijiet vojta għal kull oġġett fuq l{-}grilji tiegħna tnejn. Fejn qegħdin l{-}oġġetti tiegĘek ISSA (koordinati x'jkollok)?} \\
            }
        }
    }
     \\ \\

    \theutterance \stepcounter{utterance}  
    & & & \multicolumn{2}{p{0.3\linewidth}}{
        \cellcolor[rgb]{0.9,0.9,0.9}{
            \makecell[{{p{\linewidth}}}]{
                \texttt{\tiny{[GM$|$GM]}}
                \texttt{GĦID: F'Riga 1 tiegħi, l{-}ċelloli vojta huma: (1,1) W, (3,1) I, (4,1) T, (5,1) C, (7,1) H. L{-}ċelloli mblokkati mill{-}ħajtar huma: (2,1) '│', (6,1) '│'. Madwar l{-}oġġetti tiegħi: I f'(3,1) tista' tmur (4,1); T f'(4,1) mhux movibbli; C f'(5,1) tista' tmur (4,1) jew (7,1) mhux vojta; H f'(7,1) mhux movibbli. } \\
\\ 
\texttt{Għalhekk, il{-}movimenti possibbli f'Riga 1 huma limitati. Nista' nċaqlaq 'I' minn (3,1) għal (4,1) meta tkun vojta. } \\
\\ 
\texttt{L{-}ebda oġġett ma jista' jidħol f'(2,1) jew (6,1) għax huma ħajtar permanenti. } \\
\\ 
\texttt{Għandna bżonn niftakru għan realistiku li jikseb pożizzjonijiet vojta għal kull oġġett fuq l{-}grilji tiegħna tnejn. Fejn qegħdin l{-}oġġetti tiegĘek ISSA (koordinati x'jkollok)?} \\
            }
        }
    }
    & & \\ \\

    \theutterance \stepcounter{utterance}  
    & & & \multicolumn{2}{p{0.3\linewidth}}{
        \cellcolor[rgb]{0.9,0.9,0.9}{
            \makecell[{{p{\linewidth}}}]{
                \texttt{\tiny{[GM$|$GM]}}
                \texttt{Invalid format: tail is not empty: '} \\
\\ 
\texttt{Għalhekk, il{-}movimenti possibbli f'Riga 1 huma limitati. Nista' nċaqlaq 'I' minn (3,1) għal (4,1) meta tkun vojta. } \\
\\ 
\texttt{L{-}ebda oġġett ma jista' jidħol f'(2,1) jew (6,1) għax huma ħajtar permanenti. } \\
\\ 
\texttt{Għandna bżonn niftakru għan realistiku li jikseb pożizzjonijiet vojta għal kull oġġett fuq l{-}grilji tiegħna tnejn. Fejn qegħdin l{-}oġġetti tiegĘek ISSA (koordinati x'jkollok)?} \\
            }
        }
    }
    & & \\ \\

    \theutterance \stepcounter{utterance}  
    & & \multicolumn{4}{p{0.6\linewidth}}{
        \cellcolor[rgb]{0.9,0.9,0.9}{
            \makecell[{{p{\linewidth}}}]{
                \texttt{\tiny{[GM$\rangle$P2]}}
                \texttt{Penali: Il{-}messaġġ tiegħek m'għandux ikun fih xejn wara l{-}kmand!} \\
\\ 
\texttt{Kun żgur li r{-}risposta tiegħek fiha biss jew `GĦID: <MESSAĠĠ>` jew `ĊAQLEQ: <OĠĠETT>, (<X>, <Y>)`, u xejn ieħor!} \\
\\ 
\texttt{Intom flimkien akkumulajtu 11 minn 12 penali.} \\
\texttt{Jekk jogħġbok erġa' pprova!} \\
            }
        }
    }
    & \\ \\

    \theutterance \stepcounter{utterance}  
    & & & \multicolumn{4}{p{0.6\linewidth}}{
        \cellcolor[rgb]{0.9,0.9,0.9}{
            \makecell[{{p{\linewidth}}}]{
                \texttt{\tiny{[GM$\langle$P2]}}
                \texttt{GĦID: F'Riga 1 tiegħi, l{-}ċelloli vojta huma: (1,1) W, (3,1) I, (4,1) T, (5,1) C, (7,1) H. L{-}ċelloli mblokkati mill{-}ħajtar huma: (2,1) '│', (6,1) '│'. Madwar l{-}oġġetti tiegħi: I f'(3,1) tista' tmur (4,1); T f'(4,1) mhux movibbli; C f'(5,1) tista' tmur (4,1) jew (7,1) mhux vojta; H f'(7,1) mhux movibbli. Għalhekk, il{-}movimenti possibbli f'Riga 1 huma limitati. Nista' nċaqlaq 'I' minn (3,1) għal (4,1) meta tkun vojta. L{-}ebda oġġett ma jista' jidħol f'(2,1) jew (6,1) għax huma ħajtar permanenti. Għandna bżonn niftakru għan realistiku li jikseb pożizzjonijiet vojta għal kull oġġett fuq l{-}grilji tiegħna tnejn. Fejn qegħdin l{-}oġġetti tiegħek ISSA?} \\
            }
        }
    }
     \\ \\

    \theutterance \stepcounter{utterance}  
    & & & \multicolumn{2}{p{0.3\linewidth}}{
        \cellcolor[rgb]{0.9,0.9,0.9}{
            \makecell[{{p{\linewidth}}}]{
                \texttt{\tiny{[GM$|$GM]}}
                \texttt{GĦID: F'Riga 1 tiegħi, l{-}ċelloli vojta huma: (1,1) W, (3,1) I, (4,1) T, (5,1) C, (7,1) H. L{-}ċelloli mblokkati mill{-}ħajtar huma: (2,1) '│', (6,1) '│'. Madwar l{-}oġġetti tiegħi: I f'(3,1) tista' tmur (4,1); T f'(4,1) mhux movibbli; C f'(5,1) tista' tmur (4,1) jew (7,1) mhux vojta; H f'(7,1) mhux movibbli. Għalhekk, il{-}movimenti possibbli f'Riga 1 huma limitati. Nista' nċaqlaq 'I' minn (3,1) għal (4,1) meta tkun vojta. L{-}ebda oġġett ma jista' jidħol f'(2,1) jew (6,1) għax huma ħajtar permanenti. Għandna bżonn niftakru għan realistiku li jikseb pożizzjonijiet vojta għal kull oġġett fuq l{-}grilji tiegħna tnejn. Fejn qegħdin l{-}oġġetti tiegħek ISSA?} \\
            }
        }
    }
    & & \\ \\

    \theutterance \stepcounter{utterance}  
    & & \multicolumn{4}{p{0.6\linewidth}}{
        \cellcolor[rgb]{0.9,0.9,0.9}{
            \makecell[{{p{\linewidth}}}]{
                \texttt{\tiny{[P1$\langle$GM]}}
                \texttt{Il{-}messaġġ tiegħek ġie mgħoddi lill{-}plejer l{-}ieħor.} \\
\\ 
\\ 
\texttt{Bħalissa qed tilagħbu r{-}round 17 minn massimu ta' 20.} \\
\texttt{Intom flimkien akkumulajtu 11 minn 12 penali.} \\
\texttt{Il{-}plejer l{-}ieħor bagħat dan il{-}messaġġ:} \\
\\ 
\\ 
\texttt{"F'Riga 1 tiegħi, l{-}ċelloli vojta huma: (1,1) W, (3,1) I, (4,1) T, (5,1) C, (7,1) H. L{-}ċelloli mblokkati mill{-}ħajtar huma: (2,1) '│', (6,1) '│'. Madwar l{-}oġġetti tiegħi: I f'(3,1) tista' tmur (4,1); T f'(4,1) mhux movibbli; C f'(5,1) tista' tmur (4,1) jew (7,1) mhux vojta; H f'(7,1) mhux movibbli. Għalhekk, il{-}movimenti possibbli f'Riga 1 huma limitati. Nista' nċaqlaq 'I' minn (3,1) għal (4,1) meta tkun vojta. L{-}ebda oġġett ma jista' jidħol f'(2,1) jew (6,1) għax huma ħajtar permanenti. Għandna bżonn niftakru għan realistiku li jikseb pożizzjonijiet vojta għal kull oġġett fuq l{-}grilji tiegħna tnejn. Fejn qegħdin l{-}oġġetti tiegħek ISSA?"} \\
\\ 
\\ 
\texttt{X'inhu l{-}kmand li jmiss tiegħek?} \\
            }
        }
    }
    & \\ \\

    \theutterance \stepcounter{utterance}  
    & \multicolumn{4}{p{0.6\linewidth}}{
        \cellcolor[rgb]{0.9,0.9,0.9}{
            \makecell[{{p{\linewidth}}}]{
                \texttt{\tiny{[P1$\rangle$GM]}}
                \texttt{GĦID: L{-}oġġetti tiegħi: W(1,1), I(7,1), T(3,1), C(4,1), H(5,1). (6,1) mhix vojta (fih '│'). Il{-}pożizzjonijiet vojta f'Riga 1 tiegħi huma: (2,1) '│' mblokkata, (3,1) T, (4,1) C, (5,1) H, (6,1) '│' mblokkata, (7,1) I. L{-}ebda pożizzjoni oħra f'Riga 1 mhix vojta. Għandna bżonn niddikjaraw għan realistiku b'baze tal{-}ħajtar. X'pożizzjonijiet vojta għandek int f'kollox il{-}grilja?} \\
            }
        }
    }
    & & \\ \\

    \theutterance \stepcounter{utterance}  
    & & & \multicolumn{2}{p{0.3\linewidth}}{
        \cellcolor[rgb]{0.9,0.9,0.9}{
            \makecell[{{p{\linewidth}}}]{
                \texttt{\tiny{[GM$|$GM]}}
                \texttt{GĦID: L{-}oġġetti tiegħi: W(1,1), I(7,1), T(3,1), C(4,1), H(5,1). (6,1) mhix vojta (fih '│'). Il{-}pożizzjonijiet vojta f'Riga 1 tiegħi huma: (2,1) '│' mblokkata, (3,1) T, (4,1) C, (5,1) H, (6,1) '│' mblokkata, (7,1) I. L{-}ebda pożizzjoni oħra f'Riga 1 mhix vojta. Għandna bżonn niddikjaraw għan realistiku b'baze tal{-}ħajtar. X'pożizzjonijiet vojta għandek int f'kollox il{-}grilja?} \\
            }
        }
    }
    & & \\ \\

    \theutterance \stepcounter{utterance}  
    & & \multicolumn{4}{p{0.6\linewidth}}{
        \cellcolor[rgb]{0.9,0.9,0.9}{
            \makecell[{{p{\linewidth}}}]{
                \texttt{\tiny{[GM$\rangle$P2]}}
                \texttt{Il{-}messaġġ tiegħek ġie mgħoddi lill{-}plejer l{-}ieħor.} \\
\\ 
\\ 
\texttt{Bħalissa qed tilagħbu r{-}round 18 minn massimu ta' 20.} \\
\texttt{Intom flimkien akkumulajtu 11 minn 12 penali.} \\
\texttt{Il{-}plejer l{-}ieħor bagħat dan il{-}messaġġ:} \\
\\ 
\\ 
\texttt{"L{-}oġġetti tiegħi: W(1,1), I(7,1), T(3,1), C(4,1), H(5,1). (6,1) mhix vojta (fih '│'). Il{-}pożizzjonijiet vojta f'Riga 1 tiegħi huma: (2,1) '│' mblokkata, (3,1) T, (4,1) C, (5,1) H, (6,1) '│' mblokkata, (7,1) I. L{-}ebda pożizzjoni oħra f'Riga 1 mhix vojta. Għandna bżonn niddikjaraw għan realistiku b'baze tal{-}ħajtar. X'pożizzjonijiet vojta għandek int f'kollox il{-}grilja?"} \\
\\ 
\\ 
\texttt{X'inhu l{-}kmand li jmiss tiegħek?} \\
            }
        }
    }
    & \\ \\

    \theutterance \stepcounter{utterance}  
    & & & \multicolumn{4}{p{0.6\linewidth}}{
        \cellcolor[rgb]{0.9,0.9,0.9}{
            \makecell[{{p{\linewidth}}}]{
                \texttt{\tiny{[GM$\langle$P2]}}
                \texttt{GĦID: L{-}ċelloli vojta f'grilja tiegħi: (1,1), (3,1), (4,1), (5,1), (7,1), (2,3), (3,3), (4,3), (5,3), (6,3), (7,3), (1,4), (2,4), (3,4), (4,4), (5,4), (6,4), (7,4), (1,5){-}(7,5), (1,6){-}(7,6), (1,7){-}(7,7). Mblokkati: (2,1), (6,1), (2,2), (6,2), (4,4). L{-}oġġetti tiegħi: W(1,1), I(3,1), T(4,1), C(5,1), H(7,1). Għan realistiku: nistgħu nġibuhom lura fl{-}istess kordinati ta' qabel (W(1,1), I(3,1), T(4,1), C(5,1), H(7,1)) u inti tmexxi I tiegħek għal (3,1) jew (6,1) mhux possibbli. X'raġ?} \\
            }
        }
    }
     \\ \\

    \theutterance \stepcounter{utterance}  
    & & & \multicolumn{2}{p{0.3\linewidth}}{
        \cellcolor[rgb]{0.9,0.9,0.9}{
            \makecell[{{p{\linewidth}}}]{
                \texttt{\tiny{[GM$|$GM]}}
                \texttt{GĦID: L{-}ċelloli vojta f'grilja tiegħi: (1,1), (3,1), (4,1), (5,1), (7,1), (2,3), (3,3), (4,3), (5,3), (6,3), (7,3), (1,4), (2,4), (3,4), (4,4), (5,4), (6,4), (7,4), (1,5){-}(7,5), (1,6){-}(7,6), (1,7){-}(7,7). Mblokkati: (2,1), (6,1), (2,2), (6,2), (4,4). L{-}oġġetti tiegħi: W(1,1), I(3,1), T(4,1), C(5,1), H(7,1). Għan realistiku: nistgħu nġibuhom lura fl{-}istess kordinati ta' qabel (W(1,1), I(3,1), T(4,1), C(5,1), H(7,1)) u inti tmexxi I tiegħek għal (3,1) jew (6,1) mhux possibbli. X'raġ?} \\
            }
        }
    }
    & & \\ \\

    \theutterance \stepcounter{utterance}  
    & & \multicolumn{4}{p{0.6\linewidth}}{
        \cellcolor[rgb]{0.9,0.9,0.9}{
            \makecell[{{p{\linewidth}}}]{
                \texttt{\tiny{[P1$\langle$GM]}}
                \texttt{Il{-}messaġġ tiegħek ġie mgħoddi lill{-}plejer l{-}ieħor.} \\
\\ 
\\ 
\texttt{Bħalissa qed tilagħbu r{-}round 18 minn massimu ta' 20.} \\
\texttt{Intom flimkien akkumulajtu 11 minn 12 penali.} \\
\texttt{Il{-}plejer l{-}ieħor bagħat dan il{-}messaġġ:} \\
\\ 
\\ 
\texttt{"L{-}ċelloli vojta f'grilja tiegħi: (1,1), (3,1), (4,1), (5,1), (7,1), (2,3), (3,3), (4,3), (5,3), (6,3), (7,3), (1,4), (2,4), (3,4), (4,4), (5,4), (6,4), (7,4), (1,5){-}(7,5), (1,6){-}(7,6), (1,7){-}(7,7). Mblokkati: (2,1), (6,1), (2,2), (6,2), (4,4). L{-}oġġetti tiegħi: W(1,1), I(3,1), T(4,1), C(5,1), H(7,1). Għan realistiku: nistgħu nġibuhom lura fl{-}istess kordinati ta' qabel (W(1,1), I(3,1), T(4,1), C(5,1), H(7,1)) u inti tmexxi I tiegħek għal (3,1) jew (6,1) mhux possibbli. X'raġ?"} \\
\\ 
\\ 
\texttt{X'inhu l{-}kmand li jmiss tiegħek?} \\
            }
        }
    }
    & \\ \\

    \theutterance \stepcounter{utterance}  
    & \multicolumn{4}{p{0.6\linewidth}}{
        \cellcolor[rgb]{0.9,0.9,0.9}{
            \makecell[{{p{\linewidth}}}]{
                \texttt{\tiny{[P1$\rangle$GM]}}
                \texttt{GĦID: L{-}oġġetti tiegħi: W(1,1), I(7,1), T(3,1), C(4,1), H(5,1). Il{-}kordinati vojta f'Riga 1 tiegħi: (3,1), (4,1), (5,1) mġemmurin. (2,1), (6,1), (7,1) mblokkati. L{-}oġġetti tiegħi movibbli: I(7,1) {-}> (7,2), (6,1)mblokkat; C(4,1) {-}> (4,2); H(5,1) {-}> (5,2); T(3,1) mhux movibbli. Għan realistiku: nistgħu nqablu li l{-}oġġetti identiċi jispiċċaw f'kordinati li huma vojta fuq l{-}grilji tnejn? Minbarra r{-}Riga 1? X'inhu l{-}pool ta' kordinati komuni vojta?} \\
            }
        }
    }
    & & \\ \\

    \theutterance \stepcounter{utterance}  
    & & & \multicolumn{2}{p{0.3\linewidth}}{
        \cellcolor[rgb]{0.9,0.9,0.9}{
            \makecell[{{p{\linewidth}}}]{
                \texttt{\tiny{[GM$|$GM]}}
                \texttt{GĦID: L{-}oġġetti tiegħi: W(1,1), I(7,1), T(3,1), C(4,1), H(5,1). Il{-}kordinati vojta f'Riga 1 tiegħi: (3,1), (4,1), (5,1) mġemmurin. (2,1), (6,1), (7,1) mblokkati. L{-}oġġetti tiegħi movibbli: I(7,1) {-}> (7,2), (6,1)mblokkat; C(4,1) {-}> (4,2); H(5,1) {-}> (5,2); T(3,1) mhux movibbli. Għan realistiku: nistgħu nqablu li l{-}oġġetti identiċi jispiċċaw f'kordinati li huma vojta fuq l{-}grilji tnejn? Minbarra r{-}Riga 1? X'inhu l{-}pool ta' kordinati komuni vojta?} \\
            }
        }
    }
    & & \\ \\

    \theutterance \stepcounter{utterance}  
    & & \multicolumn{4}{p{0.6\linewidth}}{
        \cellcolor[rgb]{0.9,0.9,0.9}{
            \makecell[{{p{\linewidth}}}]{
                \texttt{\tiny{[GM$\rangle$P2]}}
                \texttt{Il{-}messaġġ tiegħek ġie mgħoddi lill{-}plejer l{-}ieħor.} \\
\\ 
\\ 
\texttt{Bħalissa qed tilagħbu r{-}round 19 minn massimu ta' 20.} \\
\texttt{Intom flimkien akkumulajtu 11 minn 12 penali.} \\
\texttt{Il{-}plejer l{-}ieħor bagħat dan il{-}messaġġ:} \\
\\ 
\\ 
\texttt{"L{-}oġġetti tiegħi: W(1,1), I(7,1), T(3,1), C(4,1), H(5,1). Il{-}kordinati vojta f'Riga 1 tiegħi: (3,1), (4,1), (5,1) mġemmurin. (2,1), (6,1), (7,1) mblokkati. L{-}oġġetti tiegħi movibbli: I(7,1) {-}> (7,2), (6,1)mblokkat; C(4,1) {-}> (4,2); H(5,1) {-}> (5,2); T(3,1) mhux movibbli. Għan realistiku: nistgħu nqablu li l{-}oġġetti identiċi jispiċċaw f'kordinati li huma vojta fuq l{-}grilji tnejn? Minbarra r{-}Riga 1? X'inhu l{-}pool ta' kordinati komuni vojta?"} \\
\\ 
\\ 
\texttt{X'inhu l{-}kmand li jmiss tiegħek?} \\
            }
        }
    }
    & \\ \\

    \theutterance \stepcounter{utterance}  
    & & & \multicolumn{4}{p{0.6\linewidth}}{
        \cellcolor[rgb]{0.9,0.9,0.9}{
            \makecell[{{p{\linewidth}}}]{
                \texttt{\tiny{[GM$\langle$P2]}}
                \texttt{GĦID: L{-}istess kordinati vojta fuq l{-}grilji tnejn huma: (1,1), (3,1), (4,1), (5,1), (7,1), (1,3){-}(7,3), (1,4){-}(7,4), (1,5){-}(7,5), (1,6){-}(7,6), (1,7){-}(7,7). Mblokkati fuq kulligru: (2,1), (6,1), (2,2), (6,2), (4,4). Għan realistiku: W f'(1,1) (waqgħa spiċċat), I f'(3,1), T f'(4,1), C f'(5,1), H f'(7,1). Dak huwa l{-}aktar realistiku għax l{-}oġġetti qegħdin hemmkemm. Int għandek I f'(7,1) {-} tista' tmexxih għal (7,2) u pastaj għal (3,1) meta tkun vojta? Jew nqablu li H jispiċċa f'(7,1) u inti tmexxi I tiegħek għal (7,1) meta tkun vojta? Wara li nmexxu l{-}oġġetti tiegħi li huma f'pożizzjoni sħiħa.} \\
            }
        }
    }
     \\ \\

    \theutterance \stepcounter{utterance}  
    & & & \multicolumn{2}{p{0.3\linewidth}}{
        \cellcolor[rgb]{0.9,0.9,0.9}{
            \makecell[{{p{\linewidth}}}]{
                \texttt{\tiny{[GM$|$GM]}}
                \texttt{GĦID: L{-}istess kordinati vojta fuq l{-}grilji tnejn huma: (1,1), (3,1), (4,1), (5,1), (7,1), (1,3){-}(7,3), (1,4){-}(7,4), (1,5){-}(7,5), (1,6){-}(7,6), (1,7){-}(7,7). Mblokkati fuq kulligru: (2,1), (6,1), (2,2), (6,2), (4,4). Għan realistiku: W f'(1,1) (waqgħa spiċċat), I f'(3,1), T f'(4,1), C f'(5,1), H f'(7,1). Dak huwa l{-}aktar realistiku għax l{-}oġġetti qegħdin hemmkemm. Int għandek I f'(7,1) {-} tista' tmexxih għal (7,2) u pastaj għal (3,1) meta tkun vojta? Jew nqablu li H jispiċċa f'(7,1) u inti tmexxi I tiegħek għal (7,1) meta tkun vojta? Wara li nmexxu l{-}oġġetti tiegħi li huma f'pożizzjoni sħiħa.} \\
            }
        }
    }
    & & \\ \\

    \theutterance \stepcounter{utterance}  
    & & \multicolumn{4}{p{0.6\linewidth}}{
        \cellcolor[rgb]{0.9,0.9,0.9}{
            \makecell[{{p{\linewidth}}}]{
                \texttt{\tiny{[P1$\langle$GM]}}
                \texttt{Il{-}messaġġ tiegħek ġie mgħoddi lill{-}plejer l{-}ieħor.} \\
\\ 
\\ 
\texttt{Bħalissa qed tilagħbu r{-}round 19 minn massimu ta' 20.} \\
\texttt{Intom flimkien akkumulajtu 11 minn 12 penali.} \\
\texttt{Il{-}plejer l{-}ieħor bagħat dan il{-}messaġġ:} \\
\\ 
\\ 
\texttt{"L{-}istess kordinati vojta fuq l{-}grilji tnejn huma: (1,1), (3,1), (4,1), (5,1), (7,1), (1,3){-}(7,3), (1,4){-}(7,4), (1,5){-}(7,5), (1,6){-}(7,6), (1,7){-}(7,7). Mblokkati fuq kulligru: (2,1), (6,1), (2,2), (6,2), (4,4). Għan realistiku: W f'(1,1) (waqgħa spiċċat), I f'(3,1), T f'(4,1), C f'(5,1), H f'(7,1). Dak huwa l{-}aktar realistiku għax l{-}oġġetti qegħdin hemmkemm. Int għandek I f'(7,1) {-} tista' tmexxih għal (7,2) u pastaj għal (3,1) meta tkun vojta? Jew nqablu li H jispiċċa f'(7,1) u inti tmexxi I tiegħek għal (7,1) meta tkun vojta? Wara li nmexxu l{-}oġġetti tiegħi li huma f'pożizzjoni sħiħa."} \\
\\ 
\\ 
\texttt{X'inhu l{-}kmand li jmiss tiegħek?} \\
            }
        }
    }
    & \\ \\

    \theutterance \stepcounter{utterance}  
    & \multicolumn{4}{p{0.6\linewidth}}{
        \cellcolor[rgb]{0.9,0.9,0.9}{
            \makecell[{{p{\linewidth}}}]{
                \texttt{\tiny{[P1$\rangle$GM]}}
                \texttt{GĦID: Moqbel. Għan: W(1,1), I(3,1), T(4,1), C(5,1), H(7,1). Int tmexxi I tiegħek minn (3,1) għal (4,1) meta tkun vojta (wara li T jmur). Jien: 1) Nmexxi C minn (4,1) għal (4,2), 2) Nmexxi T minn (3,1) għal (4,1), 3) Nmexxi I minn (7,1) għal (3,1) permezz ta' (7,2) u (7,3)... imma (7,1) mblokkata għall{-}pass ta' H? (6,1) mblokkat. I staqsi: kif nista' nmexxi I minn (7,1) għal (3,1) jekk (6,1) mblokkata? Nista' nmur (7,2){-}>(7,3){-}>...{-}>(3,3){-}>(3,2){-}>(3,1) jekk vojta. Tista' tkonfirma li (7,2), (7,3), (3,3), (3,2) huma vojta f'grilja tiegħi?} \\
            }
        }
    }
    & & \\ \\

    \theutterance \stepcounter{utterance}  
    & & & \multicolumn{2}{p{0.3\linewidth}}{
        \cellcolor[rgb]{0.9,0.9,0.9}{
            \makecell[{{p{\linewidth}}}]{
                \texttt{\tiny{[GM$|$GM]}}
                \texttt{GĦID: Moqbel. Għan: W(1,1), I(3,1), T(4,1), C(5,1), H(7,1). Int tmexxi I tiegħek minn (3,1) għal (4,1) meta tkun vojta (wara li T jmur). Jien: 1) Nmexxi C minn (4,1) għal (4,2), 2) Nmexxi T minn (3,1) għal (4,1), 3) Nmexxi I minn (7,1) għal (3,1) permezz ta' (7,2) u (7,3)... imma (7,1) mblokkata għall{-}pass ta' H? (6,1) mblokkat. I staqsi: kif nista' nmexxi I minn (7,1) għal (3,1) jekk (6,1) mblokkata? Nista' nmur (7,2){-}>(7,3){-}>...{-}>(3,3){-}>(3,2){-}>(3,1) jekk vojta. Tista' tkonfirma li (7,2), (7,3), (3,3), (3,2) huma vojta f'grilja tiegħi?} \\
            }
        }
    }
    & & \\ \\

    \theutterance \stepcounter{utterance}  
    & & \multicolumn{4}{p{0.6\linewidth}}{
        \cellcolor[rgb]{0.9,0.9,0.9}{
            \makecell[{{p{\linewidth}}}]{
                \texttt{\tiny{[GM$\rangle$P2]}}
                \texttt{Il{-}messaġġ tiegħek ġie mgħoddi lill{-}plejer l{-}ieħor.} \\
\\ 
\\ 
\texttt{Bħalissa qed tilagħbu r{-}round 20 minn massimu ta' 20.} \\
\texttt{Intom flimkien akkumulajtu 11 minn 12 penali.} \\
\texttt{Il{-}plejer l{-}ieħor bagħat dan il{-}messaġġ:} \\
\\ 
\\ 
\texttt{"Moqbel. Għan: W(1,1), I(3,1), T(4,1), C(5,1), H(7,1). Int tmexxi I tiegħek minn (3,1) għal (4,1) meta tkun vojta (wara li T jmur). Jien: 1) Nmexxi C minn (4,1) għal (4,2), 2) Nmexxi T minn (3,1) għal (4,1), 3) Nmexxi I minn (7,1) għal (3,1) permezz ta' (7,2) u (7,3)... imma (7,1) mblokkata għall{-}pass ta' H? (6,1) mblokkat. I staqsi: kif nista' nmexxi I minn (7,1) għal (3,1) jekk (6,1) mblokkata? Nista' nmur (7,2){-}>(7,3){-}>...{-}>(3,3){-}>(3,2){-}>(3,1) jekk vojta. Tista' tkonfirma li (7,2), (7,3), (3,3), (3,2) huma vojta f'grilja tiegħi?"} \\
\\ 
\\ 
\texttt{X'inhu l{-}kmand li jmiss tiegħek?} \\
            }
        }
    }
    & \\ \\

    \theutterance \stepcounter{utterance}  
    & & & \multicolumn{4}{p{0.6\linewidth}}{
        \cellcolor[rgb]{0.9,0.9,0.9}{
            \makecell[{{p{\linewidth}}}]{
                \texttt{\tiny{[GM$\langle$P2]}}
                \texttt{GĦID: Iva, (7,2), (7,3), (3,3), (3,2) huma kollha vojta f'grilja tiegħi. Il{-}pass (7,1){-}>(7,2){-}>(7,3){-}>(3,3){-}>(3,2){-}>(3,1) huwa possibbli. Plan tiegħek jixrob. Jien nkompli b'plan tiegħi: 1) Nmexxi C minn (5,1) għal (4,1) meta T timxi, 2) Nmexxi T minn (4,1) għal (5,1) jew (6,1) mhux possibbli, forsi (4,2)? Għalhekk, nmexxi: 1) T (4,1){-}>(4,2), 2) C (5,1){-}>(4,1), 3) I (3,1){-}>(5,1) mhux possibbli direkt. L{-}aqwa: Nmexxi I (3,1){-}>(4,1) meta C timxi, C (5,1){-}>(3,1)? Mhux possibbli. } \\
\\ 
\texttt{Simplifika: Għan finali: W(1,1), I(3,1), T(4,1), C(5,1), H(7,1). Situazzjoni attwali: Jien: W(1,1), I(3,1), T(4,1), C(5,1), H(7,1). Int: W(1,1), I(7,1), T(3,1), C(4,1), H(5,1).} \\
\\ 
\texttt{Passi tiegħi: 1. T (4,1){-}>(4,2) 2. C (5,1){-}>(4,1) 3. I (3,1){-}>(5,1) mhux possibbli. Għalhekk: 1. T (4,1){-}>(4,2) 2. I (3,1){-}>(4,1) 3. C (5,1){-}>(3,1) 4. H (7,1) staqsi. } \\
\\ 
\texttt{Passi tiegħek: 1. H (5,1){-}>(5,2) 2. C (4,1){-}>(5,1) 3. T (3,1){-}>(4,1) 4. I (7,1){-}>(7,2){-}>(7,3){-}>(3,3){-}>(3,2){-}>(3,1). } \\
\\ 
\texttt{Jekk inti tagħmel passi tiegħek, u jien tagħmel passi tiegħi, waslu l{-}għan. Moqbel?} \\
            }
        }
    }
     \\ \\

    \theutterance \stepcounter{utterance}  
    & & & \multicolumn{2}{p{0.3\linewidth}}{
        \cellcolor[rgb]{0.9,0.9,0.9}{
            \makecell[{{p{\linewidth}}}]{
                \texttt{\tiny{[GM$|$GM]}}
                \texttt{GĦID: Iva, (7,2), (7,3), (3,3), (3,2) huma kollha vojta f'grilja tiegħi. Il{-}pass (7,1){-}>(7,2){-}>(7,3){-}>(3,3){-}>(3,2){-}>(3,1) huwa possibbli. Plan tiegħek jixrob. Jien nkompli b'plan tiegħi: 1) Nmexxi C minn (5,1) għal (4,1) meta T timxi, 2) Nmexxi T minn (4,1) għal (5,1) jew (6,1) mhux possibbli, forsi (4,2)? Għalhekk, nmexxi: 1) T (4,1){-}>(4,2), 2) C (5,1){-}>(4,1), 3) I (3,1){-}>(5,1) mhux possibbli direkt. L{-}aqwa: Nmexxi I (3,1){-}>(4,1) meta C timxi, C (5,1){-}>(3,1)? Mhux possibbli. } \\
\\ 
\texttt{Simplifika: Għan finali: W(1,1), I(3,1), T(4,1), C(5,1), H(7,1). Situazzjoni attwali: Jien: W(1,1), I(3,1), T(4,1), C(5,1), H(7,1). Int: W(1,1), I(7,1), T(3,1), C(4,1), H(5,1).} \\
\\ 
\texttt{Passi tiegħi: 1. T (4,1){-}>(4,2) 2. C (5,1){-}>(4,1) 3. I (3,1){-}>(5,1) mhux possibbli. Għalhekk: 1. T (4,1){-}>(4,2) 2. I (3,1){-}>(4,1) 3. C (5,1){-}>(3,1) 4. H (7,1) staqsi. } \\
\\ 
\texttt{Passi tiegħek: 1. H (5,1){-}>(5,2) 2. C (4,1){-}>(5,1) 3. T (3,1){-}>(4,1) 4. I (7,1){-}>(7,2){-}>(7,3){-}>(3,3){-}>(3,2){-}>(3,1). } \\
\\ 
\texttt{Jekk inti tagħmel passi tiegħek, u jien tagħmel passi tiegħi, waslu l{-}għan. Moqbel?} \\
            }
        }
    }
    & & \\ \\

    \theutterance \stepcounter{utterance}  
    & & & \multicolumn{2}{p{0.3\linewidth}}{
        \cellcolor[rgb]{0.9,0.9,0.9}{
            \makecell[{{p{\linewidth}}}]{
                \texttt{\tiny{[GM$|$GM]}}
                \texttt{Invalid format: tail is not empty: '} \\
\\ 
\texttt{Simplifika: Għan finali: W(1,1), I(3,1), T(4,1), C(5,1), H(7,1). Situazzjoni attwali: Jien: W(1,1), I(3,1), T(4,1), C(5,1), H(7,1). Int: W(1,1), I(7,1), T(3,1), C(4,1), H(5,1).} \\
\\ 
\texttt{Passi tiegħi: 1. T (4,1){-}>(4,2) 2. C (5,1){-}>(4,1) 3. I (3,1){-}>(5,1) mhux possibbli. Għalhekk: 1. T (4,1){-}>(4,2) 2. I (3,1){-}>(4,1) 3. C (5,1){-}>(3,1) 4. H (7,1) staqsi. } \\
\\ 
\texttt{Passi tiegħek: 1. H (5,1){-}>(5,2) 2. C (4,1){-}>(5,1) 3. T (3,1){-}>(4,1) 4. I (7,1){-}>(7,2){-}>(7,3){-}>(3,3){-}>(3,2){-}>(3,1). } \\
\\ 
\texttt{Jekk inti tagħmel passi tiegħek, u jien tagħmel passi tiegħi, waslu l{-}għan. Moqbel?} \\
            }
        }
    }
    & & \\ \\

    \theutterance \stepcounter{utterance}  
    & & \multicolumn{4}{p{0.6\linewidth}}{
        \cellcolor[rgb]{0.9,0.9,0.9}{
            \makecell[{{p{\linewidth}}}]{
                \texttt{\tiny{[GM$\rangle$P2]}}
                \texttt{Penali: Il{-}messaġġ tiegħek m'għandux ikun fih xejn wara l{-}kmand!} \\
\\ 
\texttt{Kun żgur li r{-}risposta tiegħek fiha biss jew `GĦID: <MESSAĠĠ>` jew `ĊAQLEQ: <OĠĠETT>, (<X>, <Y>)`, u xejn ieħor!} \\
\\ 
\texttt{Intom flimkien akkumulajtu 12 minn 12 penali.} \\
\texttt{Jekk jogħġbok erġa' pprova!} \\
            }
        }
    }
    & \\ \\

    \theutterance \stepcounter{utterance}  
    & & & \multicolumn{4}{p{0.6\linewidth}}{
        \cellcolor[rgb]{0.9,0.9,0.9}{
            \makecell[{{p{\linewidth}}}]{
                \texttt{\tiny{[GM$\langle$P2]}}
                \texttt{GĦID: Iva, (7,2), (7,3), (3,3), (3,2) huma vojta. Plan tiegħek jixrob. Għan: W(1,1), I(3,1), T(4,1), C(5,1), H(7,1). Passi tiegħi: 1) T(4,1){-}>(4,2), 2) I(3,1){-}>(4,1), 3) C(5,1){-}>(3,1). Passi tiegħek: 1) H(5,1){-}>(5,2), 2) C(4,1){-}>(5,1), 3) T(3,1){-}>(4,1), 4) I(7,1){-}>(7,2){-}>(7,3){-}>(3,3){-}>(3,2){-}>(3,1). Moqbel?} \\
            }
        }
    }
     \\ \\

    \theutterance \stepcounter{utterance}  
    & & & \multicolumn{2}{p{0.3\linewidth}}{
        \cellcolor[rgb]{0.9,0.9,0.9}{
            \makecell[{{p{\linewidth}}}]{
                \texttt{\tiny{[GM$|$GM]}}
                \texttt{GĦID: Iva, (7,2), (7,3), (3,3), (3,2) huma vojta. Plan tiegħek jixrob. Għan: W(1,1), I(3,1), T(4,1), C(5,1), H(7,1). Passi tiegħi: 1) T(4,1){-}>(4,2), 2) I(3,1){-}>(4,1), 3) C(5,1){-}>(3,1). Passi tiegħek: 1) H(5,1){-}>(5,2), 2) C(4,1){-}>(5,1), 3) T(3,1){-}>(4,1), 4) I(7,1){-}>(7,2){-}>(7,3){-}>(3,3){-}>(3,2){-}>(3,1). Moqbel?} \\
            }
        }
    }
    & & \\ \\

    \theutterance \stepcounter{utterance}  
    & & & \multicolumn{2}{p{0.3\linewidth}}{
        \cellcolor[rgb]{0.9,0.9,0.9}{
            \makecell[{{p{\linewidth}}}]{
                \texttt{\tiny{[GM$|$GM]}}
                \texttt{Maximum number of rounds reached} \\
            }
        }
    }
    & & \\ \\

    \theutterance \stepcounter{utterance}  
    & & & \multicolumn{2}{p{0.3\linewidth}}{
        \cellcolor[rgb]{0.9,0.9,0.9}{
            \makecell[{{p{\linewidth}}}]{
                \texttt{\tiny{[GM$|$GM]}}
                \texttt{Initial states:} \\
\texttt{```} \\
\texttt{Player 1:} \\
\texttt{ 1234567} \\
\texttt{╔═╤═══╤═╗} \\
\texttt{║◌│◌◌◌│◌║ 1} \\
\texttt{╟─┼──┬┘◌║ 2} \\
\texttt{║◌│◌◌│◌◌║ 3} \\
\texttt{╟─┴──┘H◌║ 4} \\
\texttt{║◌◌◌◌◌I◌║ 5} \\
\texttt{║C◌◌◌◌◌◌║ 6} \\
\texttt{║◌W◌◌T◌◌║ 7} \\
\texttt{╚═══════╝} \\
\\ 
\\ 
\texttt{Player 2:} \\
\texttt{ 1234567} \\
\texttt{╔═╤═══╤═╗} \\
\texttt{║◌│◌◌◌│H║ 1} \\
\texttt{╟─┼──┬┘◌║ 2} \\
\texttt{║◌│◌◌│◌C║ 3} \\
\texttt{╟─┴──┘◌T║ 4} \\
\texttt{║◌◌◌◌◌◌◌║ 5} \\
\texttt{║◌◌◌I◌◌◌║ 6} \\
\texttt{║◌◌◌◌W◌◌║ 7} \\
\texttt{╚═══════╝} \\
\\ 
\texttt{```} \\
            }
        }
    }
    & & \\ \\

    \theutterance \stepcounter{utterance}  
    & & & \multicolumn{2}{p{0.3\linewidth}}{
        \cellcolor[rgb]{0.9,0.9,0.9}{
            \makecell[{{p{\linewidth}}}]{
                \texttt{\tiny{[GM$|$GM]}}
                \texttt{End states:} \\
\texttt{```} \\
\texttt{Player 1:} \\
\texttt{ 1234567} \\
\texttt{╔═╤═══╤═╗} \\
\texttt{║W│TCH│I║ 1} \\
\texttt{╟─┼──┬┘◌║ 2} \\
\texttt{║◌│◌◌│◌◌║ 3} \\
\texttt{╟─┴──┘◌◌║ 4} \\
\texttt{║◌◌◌◌◌◌◌║ 5} \\
\texttt{║◌◌◌◌◌◌◌║ 6} \\
\texttt{║◌◌◌◌◌◌◌║ 7} \\
\texttt{╚═══════╝} \\
\\ 
\\ 
\texttt{Player 2:} \\
\texttt{ 1234567} \\
\texttt{╔═╤═══╤═╗} \\
\texttt{║W│ITC│H║ 1} \\
\texttt{╟─┼──┬┘◌║ 2} \\
\texttt{║◌│◌◌│◌◌║ 3} \\
\texttt{╟─┴──┘◌◌║ 4} \\
\texttt{║◌◌◌◌◌◌◌║ 5} \\
\texttt{║◌◌◌◌◌◌◌║ 6} \\
\texttt{║◌◌◌◌◌◌◌║ 7} \\
\texttt{╚═══════╝} \\
\\ 
\texttt{```} \\
            }
        }
    }
    & & \\ \\

    \theutterance \stepcounter{utterance}  
    & & & \multicolumn{2}{p{0.3\linewidth}}{
        \cellcolor[rgb]{0.9,0.9,0.9}{
            \makecell[{{p{\linewidth}}}]{
                \texttt{\tiny{[GM$|$GM]}}
                \texttt{* success: True} \\
\texttt{* lose: False} \\
\texttt{* aborted: False} \\
\texttt{{-}{-}{-}{-}{-}{-}{-}} \\
\texttt{* Shifts: 4.00} \\
\texttt{* Max Shifts: 8.00} \\
\texttt{* Min Shifts: 4.00} \\
\texttt{* End Distance Sum: 8.00} \\
\texttt{* Init Distance Sum: 18.71} \\
\texttt{* Expected Distance Sum: 20.95} \\
\texttt{* Penalties: 12.00} \\
\texttt{* Max Penalties: 12.00} \\
\texttt{* Rounds: 20.00} \\
\texttt{* Max Rounds: 20.00} \\
\texttt{* Object Count: 5.00} \\
            }
        }
    }
    & & \\ \\

\end{supertabular}
}

\end{document}
